\newpage
\section{Overview of the experiments}

The experimental analysis is divided into two sections, the concerning uniform and non-uniform query costs. In the first section we measured the running time of the \FRankingBasedDT algorithm from Section~\ref{sec:uniform_cost_algorithm} and the depth of the decision trees produced by it. The second secition was concerned with the non-uniform cost model and included the evaluation of the $O\br{\log n/\log\log n}$-approximation algorithm from Section~\ref{sec:log n/log log n approximation}, the $O\br{\sqrt{\log n}}$-approximation algorithm from Section~\ref{sec:sqrt log n approximation} based on the QPTAS and the $O\br{\log\log n}$-approximation algorithm from Section~\ref{sec:log log n approximation}. We compared the quality of the solutions produced by these algorithms on various classes of trees and cost functions and the computation time required to obtain these results. Additionally, for small input sizes we menaged to calculate the optimal solutions using the exact exponential time algorithm. This allowed us to compare the approximation ratios of the algorithms with respect to the optimal solutions.

We have also employed few heurstics speed-ups particularly for the QPTAS. Firstly, instead of using $\epsilon=1$ we set it to $\epsilon=59$ which reduces most of the computational complexity of each QPTAS call. For both the $O\br{\sqrt{\log n}}$-approximation algorithm and the $O\br{\log\log n}$-approximation algorithm this does not impact the asymptotical guarantee of the algorithms. Secondly, a careful reading of the proof of the QPTAS reveals that the depth of the considered boxed decision trees can be limited to $\rank\br{T}$ where $\rank\br{T}$ is the number of labels used by the optimal ranking of $T$. This reduces the running time of the QPTAS to $n^{O\br{\log \rank\br{T}/\epsilon^2}}$ which in practice is usually significantly better.


\subsection*{Tree Generation Methods}

We used three different methods of generating random trees for our experiments: random trees, trees with bounded degree and trees with bounded diameter. 
Table~\ref{tab:tree_methods} summarizes all tree generation methods used in our experiments. We decided to measure the performance of our algorithms on relatively small instances, 10 to 100 vertices per instance, since the complex dynamic programming subroutines of our algorithm require significant computation time even for such sizes.

\begin{table}[H]
\centering
\caption{Tree generation methods and their properties}
\label{tab:tree_methods}
\small
\begin{tabular}{|l|l|p{7cm}|}
\hline
\textbf{Method} & \textbf{Parameters} & \textbf{Description} \\
\hline
Random & $n$ & Uniformly random labeled tree generated using Prüfer sequence \\
\hline
Bounded Degree & $n, \Delta$ & Tree with maximum degree constraint $\Delta$ (each vertex has degree $\leq \Delta$) \\
\hline
Bounded Diameter & $n, D$ & Tree with maximum diameter constraint $D$ (maximum distance between any two vertices $\leq D$) \\
\hline
\end{tabular}
\end{table}

\subsubsection*{Detailed Generation Algorithms}

\paragraph{Random Trees.} Generated using NetworkX's \texttt{random\_labeled\_tree} function, which employs Prüfer sequences to sample uniformly from all labeled trees on $n$ vertices. A random vertex is chosen as root, and the tree is oriented via BFS traversal from this root.

\paragraph{Bounded Degree Trees.} Constructed incrementally by maintaining a list of available parent nodes with fewer than $\Delta-1$ children. Each new vertex is attached to a uniformly random available parent. When a parent reaches its degree limit, it is removed from the available list. This ensures every vertex has total degree at most $\Delta$.

\paragraph{Bounded Diameter Trees.} For diameter bound $D$, vertices are distributed across $h = \lfloor D/2 \rfloor + 1$ levels. The root occupies level 0, and remaining vertices are randomly assigned to levels $1, \ldots, h-1$. Each non-root vertex connects to a random parent in the previous level, ensuring the diameter does not exceed $D$.

\subsubsection*{Cost Function Generation}

Table~\ref{tab:distributions} lists all probability distributions used for vertex weights ($w$) and edge costs ($c$).

\begin{table}[H]
\centering
\caption{Probability distributions for edge costs}
\label{tab:distributions}
\begin{tabular}{|l|l|p{7cm}|}
\hline
\textbf{Distribution} & \textbf{Formula} & \textbf{Properties} \\
\hline
Uniform & $\mathcal{U}(0,1)$ & All values equally likely in $[0,1]$, mean $\mathbb{E}[X]=0.5$ \\
\hline
Binomial & $\text{Binomial}(n, p)$ & Discrete distribution with parameters $n=10$, $p=0.5$, mean $\mathbb{E}[X]=np=5$ \\
\hline
Exponential & $\mathcal{E}(\lambda=1)$ & Memoryless, continuous distribution with rate $\lambda=1$, mean $\mathbb{E}[X]=1$ \\
\hline
\end{tabular}
\end{table}

\subsubsection*{Extended Three-Field Notation for Experiments}

To systematically describe our experimental configurations, we extend the theoretical three-field notation $\alpha || \beta || \gamma$ to include specific tree generation methods and distribution types:

\textbf{Field $\alpha$ (Search Space):} In experiments, we specify the specific tree generation method using subscripts. For this study, we focus on four representative tree classes:
\begin{itemize}
    \item $T_R$ - Random trees (uniformly random labeled trees via Prüfer sequence)
    \item $T_{\Delta}$ - Bounded degree trees (maximum degree constraint $\Delta$)
    \item $T_D$ - Bounded diameter trees (maximum diameter constraint $D$)
    \item $T_{bal}$ - Balanced trees (balanced $m$-ary trees with minimal height)
\end{itemize}

\textbf{Field $\beta$ (Query Characteristics):} All experiments use vertex queries ($V$). For worst-case optimization, we consider four cost scenarios:
\begin{itemize}
    \item $V$ - vertex queries with uniform unit costs ($c(v) = 1$ for all $v$)
    \item $V, c = \mathcal{U}$ - vertex queries with costs from uniform distribution $\mathcal{U}(0,1)$
    \item $V, c = \mathcal{B}$ - vertex queries with costs from binomial distribution $\text{Binomial}(n=10, p=0.5)$
    \item $V, c = \mathcal{E}$ - vertex queries with costs from exponential distribution $\mathcal{E}(\lambda=1)$
\end{itemize}

\subsubsection*{Experiment Coverage Matrix}

Table~\ref{tab:experiment_matrix} shows which combinations of tree types and problem variants are tested. Each cell contains a reference to the corresponding subsection with detailed results.

\begin{table}[H]
\centering
\caption{Experiment coverage: tree generation methods vs. problem variants}
\label{tab:experiment_matrix}
\scriptsize
\begin{tabular}{|l|c|c|}
\hline
\textbf{Tree Method} & $T || V || \COST_{max}$ & $T || V, c || \COST_{max}$ \\
\hline
Random & \ref{sec:random_trees} & \ref{sec:random_c} \\
\hline
Bounded Degree & \ref{sec:bounded_degree} & \ref{sec:bounded_degree_c} \\
\hline
Bounded Diameter & \ref{sec:bounded_diameter} & \ref{sec:bounded_diameter_c} \\
\hline
Balanced & \ref{sec:balanced_trees} & \ref{sec:balanced_c} \\
\hline
\end{tabular}
\end{table}

\newpage
\section{Trees with Uniform Costs (Worst Case)}

This section presents experimental results for the basic tree search problem with uniform query costs. The objective is to construct a decision tree that minimizes the worst-case number of queries needed to identify any target vertex.

\subsection{Random Trees ($T_R || V || \COST$)}
\label{sec:random_trees}

Uniformly random labeled trees generated using the standard random tree generation algorithm (Prüfer sequences). In the worst-case uniform cost model, all queries have cost 1, and we measure the maximum depth of the decision tree.

\begin{table}[H]
\centering
\caption{Performance on random trees with uniform query costs}
\label{tab:random_uniform}
\begin{tabular}{|c|c|c|c|c|}
\hline
$n$ & Avg Depth & Max Depth & Avg Time (s) & Max Time (s) \\
\hline
\csvreader[
    head to column names,
    filter expr={
        test{\ifnumequal{\thecsvinputline}{1}} or
        test{\ifnumequal{\Size}{10}} or
        test{\ifnumequal{\Size}{20}} or
        test{\ifnumequal{\Size}{30}} or
        test{\ifnumequal{\Size}{40}} or
        test{\ifnumequal{\Size}{50}} or
        test{\ifnumequal{\Size}{60}} or
        test{\ifnumequal{\Size}{70}} or
        test{\ifnumequal{\Size}{80}} or
        test{\ifnumequal{\Size}{90}} or
        test{\ifnumequal{\Size}{100}}
    }
]{data/analysis_basic_ranking_based_dt_results.csv}{Size=\Size,AvgCPUTime_s=\AvgCPUTimes,MaxCPUTime_s=\MaxCPUTimes,AvgDepth=\AvgDepth,MaxDepth=\MaxDepth}{%
\Size & \AvgDepth & \MaxDepth & \AvgCPUTimes & \MaxCPUTimes \\
}
\hline
\end{tabular}
\end{table}

\begin{figure}[H]
\centering
\begin{tikzpicture}
\begin{axis}[
    xlabel={Tree size ($n$)},
    ylabel={Decision tree depth},
    width=0.48\textwidth,
    height=6cm,
    grid=major,
    legend pos=north west,
    mark size=1.5pt
]
\addplot[blue, mark=*, thick] table[x=Size, y=AvgDepth, col sep=comma] {data/analysis_basic_ranking_based_dt_results.csv};
\addplot[red, mark=square*, thick] table[x=Size, y=MaxDepth, col sep=comma] {data/analysis_basic_ranking_based_dt_results.csv};
\legend{Avg Depth, Max Depth}
\end{axis}
\end{tikzpicture}
\hfill
\begin{tikzpicture}
\begin{axis}[
    xlabel={Tree size ($n$)},
    ylabel={Runtime (s)},
    width=0.48\textwidth,
    height=6cm,
    grid=major,
    legend pos=north west,
    mark size=1.5pt
]
\addplot[blue, mark=*, thick] table[x=Size, y=AvgCPUTime_s, col sep=comma] {data/analysis_basic_ranking_based_dt_results.csv};
\addplot[red, mark=square*, thick] table[x=Size, y=MaxCPUTime_s, col sep=comma] {data/analysis_basic_ranking_based_dt_results.csv};
\legend{Avg Time, Max Time}
\end{axis}
\end{tikzpicture}
\caption{Performance metrics for random trees (uniform query costs)}
\label{fig:random_uniform}
\end{figure}

The average depth grows approximately as $\log_2 n$, consistent with theoretical bounds. Centroid decomposition achieves near-optimal decision tree height.

\subsection{Bounded Degree Trees ($T_{\Delta} || V || \COST$)}
\label{sec:bounded_degree}
\label{sec:bounded_degree}

Results for trees with maximum degree constraint $\Delta$. In the worst-case scenario, we measure the maximum depth of decision trees for different degree bounds.

\subsubsection{Degree 3}

\begin{table}[H]
\centering
\caption{Trees with $\Delta = 3$ (worst case depth)}
\label{tab:degree3_uniform}
\begin{tabular}{|c|c|c|c|c|}
\hline
$n$ & Avg Depth & Max Depth & Avg Time (ms) & Std Dev \\
\hline
% Data from: bounded_degree_3_uniform.csv
\csvreader[head to column names]{data/bounded_degree_3_uniform.csv}{}{%
\n & \avgdepth & \maxdepth & \avgtime & \stddev \\
}
\hline
\end{tabular}
\end{table}
\begin{figure}[H]
\centering
\begin{tikzpicture}
\begin{axis}[
    xlabel={Tree size ($n$)},
    ylabel={Average depth},
    width=0.48\textwidth,
    height=6cm,
    grid=major,
    legend pos=north west,
    mark size=2pt
]
\addplot[blue, mark=*, thick] table[x=n, y=avgdepth, col sep=comma] {data/bounded_degree_3_uniform.csv};
\legend{Avg Depth}
\end{axis}
\end{tikzpicture}
\hfill
\begin{tikzpicture}
\begin{axis}[
    xlabel={Tree size ($n$)},
    ylabel={Runtime (ms)},
    width=0.48\textwidth,
    height=6cm,
    grid=major,
    legend pos=north west,
    mark size=2pt
]
\addplot[red, mark=square*, thick] table[x=n, y=avgtime, col sep=comma] {data/bounded_degree_3_uniform.csv};
\legend{Avg Time}
\end{axis}
\end{tikzpicture}
\caption{Performance metrics: Trees with $\Delta = 3$}
\label{fig:degree3_uniform}
\end{figure}


\subsubsection{Degree 5}

\begin{table}[H]
\centering
\caption{Trees with $\Delta = 5$ (worst case depth)}
\label{tab:degree5_uniform}
\begin{tabular}{|c|c|c|c|c|}
\hline
$n$ & Avg Depth & Max Depth & Avg Time (ms) & Std Dev \\
\hline
% Data from: bounded_degree_5_uniform.csv
\csvreader[head to column names]{data/bounded_degree_5_uniform.csv}{}{%
\n & \avgdepth & \maxdepth & \avgtime & \stddev \\
}
\hline
\end{tabular}
\end{table}
\begin{figure}[H]
\centering
\begin{tikzpicture}
\begin{axis}[
    xlabel={Tree size ($n$)},
    ylabel={Average depth},
    width=0.48\textwidth,
    height=6cm,
    grid=major,
    legend pos=north west,
    mark size=2pt
]
\addplot[blue, mark=*, thick] table[x=n, y=avgdepth, col sep=comma] {data/bounded_degree_5_uniform.csv};
\legend{Avg Depth}
\end{axis}
\end{tikzpicture}
\hfill
\begin{tikzpicture}
\begin{axis}[
    xlabel={Tree size ($n$)},
    ylabel={Runtime (ms)},
    width=0.48\textwidth,
    height=6cm,
    grid=major,
    legend pos=north west,
    mark size=2pt
]
\addplot[red, mark=square*, thick] table[x=n, y=avgtime, col sep=comma] {data/bounded_degree_5_uniform.csv};
\legend{Avg Time}
\end{axis}
\end{tikzpicture}
\caption{Performance metrics: Trees with $\Delta = 5$}
\label{fig:degree5_uniform}
\end{figure}


\subsubsection{Degree 10}

\begin{table}[H]
\centering
\caption{Trees with $\Delta = 10$ (worst case depth)}
\label{tab:degree10_uniform}
\begin{tabular}{|c|c|c|c|c|}
\hline
$n$ & Avg Depth & Max Depth & Avg Time (ms) & Std Dev \\
\hline
% Data from: bounded_degree_10_uniform.csv
\csvreader[head to column names]{data/bounded_degree_10_uniform.csv}{}{%
\n & \avgdepth & \maxdepth & \avgtime & \stddev \\
}
\hline
\end{tabular}
\end{table}
\begin{figure}[H]
\centering
\begin{tikzpicture}
\begin{axis}[
    xlabel={Tree size ($n$)},
    ylabel={Average depth},
    width=0.48\textwidth,
    height=6cm,
    grid=major,
    legend pos=north west,
    mark size=2pt
]
\addplot[blue, mark=*, thick] table[x=n, y=avgdepth, col sep=comma] {data/bounded_degree_10_uniform.csv};
\legend{Avg Depth}
\end{axis}
\end{tikzpicture}
\hfill
\begin{tikzpicture}
\begin{axis}[
    xlabel={Tree size ($n$)},
    ylabel={Runtime (ms)},
    width=0.48\textwidth,
    height=6cm,
    grid=major,
    legend pos=north west,
    mark size=2pt
]
\addplot[red, mark=square*, thick] table[x=n, y=avgtime, col sep=comma] {data/bounded_degree_10_uniform.csv};
\legend{Avg Time}
\end{axis}
\end{tikzpicture}
\caption{Performance metrics: Trees with $\Delta = 10$}
\label{fig:degree10_uniform}
\end{figure}


\subsubsection{Comparison Across Degree Bounds}

\begin{table}[H]
\centering
\caption{Performance comparison for different degree bounds ($n = 100$)}
\label{tab:degree_comparison}
\begin{tabular}{|c|c|c|c|c|}
\hline
$\Delta$ & Avg Depth & Max Depth & Avg Time (ms) & Std Dev \\
\hline
% Data from: bounded_degree_comparison.csv
\csvreader[head to column names]{data/bounded_degree_comparison.csv}{}{%
\delta & \avgdepth & \maxdepth & \avgtime & \stddev \\
}
\hline
\end{tabular}
\end{table}
\begin{figure}[H]
\centering
\begin{tikzpicture}
\begin{axis}[
    xlabel={Tree size ($n$)},
    ylabel={Average depth},
    width=0.48\textwidth,
    height=6cm,
    grid=major,
    legend pos=north west,
    mark size=2pt
]
\addplot[blue, mark=*, thick] table[x=n, y=avgdepth, col sep=comma] {data/bounded_degree_comparison.csv};
\legend{Avg Depth}
\end{axis}
\end{tikzpicture}
\hfill
\begin{tikzpicture}
\begin{axis}[
    xlabel={Tree size ($n$)},
    ylabel={Runtime (ms)},
    width=0.48\textwidth,
    height=6cm,
    grid=major,
    legend pos=north west,
    mark size=2pt
]
\addplot[red, mark=square*, thick] table[x=n, y=avgtime, col sep=comma] {data/bounded_degree_comparison.csv};
\legend{Avg Time}
\end{axis}
\end{tikzpicture}
\caption{Performance metrics: Performance comparison for different degree bounds ($n = 100$)}
\label{fig:degree_comparison}
\end{figure}


\textbf{Observations:}
\begin{itemize}
    \item Lower degree bounds lead to deeper decision trees due to structural constraints
    \item As $\Delta$ increases, maximum depth decreases as trees can be more balanced
    \item Degree constraint forces more path-like structures, increasing variance
    \item Runtime remains efficient across all degree bounds
    \item Structural properties alone determine worst-case depth
\end{itemize}

\subsection{Bounded Diameter Trees ($T_D || V || \COST$)}
\label{sec:bounded_diameter}

Results for trees with diameter constraint $D$. The diameter constraint naturally limits the maximum depth of decision trees.

\subsubsection{Diameter 5}

\begin{table}[H]
\centering
\caption{Trees with $D = 5$ (worst case depth)}
\label{tab:diam5_uniform}
\begin{tabular}{|c|c|c|c|c|}
\hline
$n$ & Avg Depth & Max Depth & Avg Time (ms) & Std Dev \\
\hline
% Data from: bounded_diameter_5_uniform.csv
\csvreader[head to column names]{data/bounded_diameter_5_uniform.csv}{}{%
\n & \avgdepth & \maxdepth & \avgtime & \stddev \\
}
\hline
\end{tabular}
\end{table}
\begin{figure}[H]
\centering
\begin{tikzpicture}
\begin{axis}[
    xlabel={Tree size ($n$)},
    ylabel={Average depth},
    width=0.48\textwidth,
    height=6cm,
    grid=major,
    legend pos=north west,
    mark size=2pt
]
\addplot[blue, mark=*, thick] table[x=n, y=avgdepth, col sep=comma] {data/bounded_diameter_5_uniform.csv};
\legend{Avg Depth}
\end{axis}
\end{tikzpicture}
\hfill
\begin{tikzpicture}
\begin{axis}[
    xlabel={Tree size ($n$)},
    ylabel={Runtime (ms)},
    width=0.48\textwidth,
    height=6cm,
    grid=major,
    legend pos=north west,
    mark size=2pt
]
\addplot[red, mark=square*, thick] table[x=n, y=avgtime, col sep=comma] {data/bounded_diameter_5_uniform.csv};
\legend{Avg Time}
\end{axis}
\end{tikzpicture}
\caption{Performance metrics: Trees with $D = 5$}
\label{fig:diam5_uniform}
\end{figure}


\subsubsection{Diameter 10}

\begin{table}[H]
\centering
\caption{Trees with $D = 10$ (worst case depth)}
\label{tab:diam10_uniform}
\begin{tabular}{|c|c|c|c|c|}
\hline
$n$ & Avg Depth & Max Depth & Avg Time (ms) & Std Dev \\
\hline
% Data from: bounded_diameter_10_uniform.csv
\csvreader[head to column names]{data/bounded_diameter_10_uniform.csv}{}{%
\n & \avgdepth & \maxdepth & \avgtime & \stddev \\
}
\hline
\end{tabular}
\end{table}
\begin{figure}[H]
\centering
\begin{tikzpicture}
\begin{axis}[
    xlabel={Tree size ($n$)},
    ylabel={Average depth},
    width=0.48\textwidth,
    height=6cm,
    grid=major,
    legend pos=north west,
    mark size=2pt
]
\addplot[blue, mark=*, thick] table[x=n, y=avgdepth, col sep=comma] {data/bounded_diameter_10_uniform.csv};
\legend{Avg Depth}
\end{axis}
\end{tikzpicture}
\hfill
\begin{tikzpicture}
\begin{axis}[
    xlabel={Tree size ($n$)},
    ylabel={Runtime (ms)},
    width=0.48\textwidth,
    height=6cm,
    grid=major,
    legend pos=north west,
    mark size=2pt
]
\addplot[red, mark=square*, thick] table[x=n, y=avgtime, col sep=comma] {data/bounded_diameter_10_uniform.csv};
\legend{Avg Time}
\end{axis}
\end{tikzpicture}
\caption{Performance metrics: Trees with $D = 10$}
\label{fig:diam10_uniform}
\end{figure}


\subsubsection{Diameter 20}

\begin{table}[H]
\centering
\caption{Trees with $D = 20$ (worst case depth)}
\label{tab:diam20_uniform}
\begin{tabular}{|c|c|c|c|c|}
\hline
$n$ & Avg Depth & Max Depth & Avg Time (ms) & Std Dev \\
\hline
% Data from: bounded_diameter_20_uniform.csv
\csvreader[head to column names]{data/bounded_diameter_20_uniform.csv}{}{%
\n & \avgdepth & \maxdepth & \avgtime & \stddev \\
}
\hline
\end{tabular}
\end{table}
\begin{figure}[H]
\centering
\begin{tikzpicture}
\begin{axis}[
    xlabel={Tree size ($n$)},
    ylabel={Average depth},
    width=0.48\textwidth,
    height=6cm,
    grid=major,
    legend pos=north west,
    mark size=2pt
]
\addplot[blue, mark=*, thick] table[x=n, y=avgdepth, col sep=comma] {data/bounded_diameter_20_uniform.csv};
\legend{Avg Depth}
\end{axis}
\end{tikzpicture}
\hfill
\begin{tikzpicture}
\begin{axis}[
    xlabel={Tree size ($n$)},
    ylabel={Runtime (ms)},
    width=0.48\textwidth,
    height=6cm,
    grid=major,
    legend pos=north west,
    mark size=2pt
]
\addplot[red, mark=square*, thick] table[x=n, y=avgtime, col sep=comma] {data/bounded_diameter_20_uniform.csv};
\legend{Avg Time}
\end{axis}
\end{tikzpicture}
\caption{Performance metrics: Trees with $D = 20$}
\label{fig:diam20_uniform}
\end{figure}


\subsubsection{Comparison Across Diameter Bounds}

\begin{table}[H]
\centering
\caption{Performance comparison for different diameter bounds ($n = 100$)}
\label{tab:diameter_comparison}
\begin{tabular}{|c|c|c|c|c|}
\hline
$D$ & Avg Depth & Max Depth & Avg Time (ms) & Std Dev \\
\hline
% Data from: bounded_diameter_comparison.csv
\csvreader[head to column names]{data/bounded_diameter_comparison.csv}{}{%
\diameter & \avgdepth & \maxdepth & \avgtime & \stddev \\
}
\hline
\end{tabular}
\end{table}
\begin{figure}[H]
\centering
\begin{tikzpicture}
\begin{axis}[
    xlabel={Tree size ($n$)},
    ylabel={Average depth},
    width=0.48\textwidth,
    height=6cm,
    grid=major,
    legend pos=north west,
    mark size=2pt
]
\addplot[blue, mark=*, thick] table[x=n, y=avgdepth, col sep=comma] {data/bounded_diameter_comparison.csv};
\legend{Avg Depth}
\end{axis}
\end{tikzpicture}
\hfill
\begin{tikzpicture}
\begin{axis}[
    xlabel={Tree size ($n$)},
    ylabel={Runtime (ms)},
    width=0.48\textwidth,
    height=6cm,
    grid=major,
    legend pos=north west,
    mark size=2pt
]
\addplot[red, mark=square*, thick] table[x=n, y=avgtime, col sep=comma] {data/bounded_diameter_comparison.csv};
\legend{Avg Time}
\end{axis}
\end{tikzpicture}
\caption{Performance metrics: Performance comparison for different diameter bounds ($n = 100$)}
\label{fig:diameter_comparison}
\end{figure}


\textbf{Observations:}
\begin{itemize}
    \item Diameter constraints naturally limit decision tree depth
    \item Smaller diameter bounds force more balanced tree structures
    \item Trees with small $D$ tend toward star-like configurations
    \item Performance remains consistent across diameter variations
    \item Trade-off between tree diameter and search efficiency
\end{itemize}

% \newpage
\section{Trees with Non-Uniform Weights (Average Case)}

This section presents experimental results for algorithms solving the average-case tree search problem with non-uniform vertex weights. We evaluate three main approaches:
\begin{itemize}
    \item Greedy algorithm (2-approximation)
    \item PTAS (Polynomial-Time Approximation Scheme)
    \item FPTAS (Fully Polynomial-Time Approximation Scheme)
\end{itemize}

\input{chapters/experimental_results/trees_average_case_non_uniform_weights/random_w.tex}
\input{chapters/experimental_results/trees_average_case_non_uniform_weights/bounded_degree_w.tex}
\subsection{Bounded Diameter Trees with Non-Uniform Weights ($T_D || V, w || \sum w_i C_i$)}
\label{sec:bounded_diameter_w}

Results for bounded diameter trees with non-uniform vertex weights.

\begin{table}[H]
\centering
\caption{Bounded diameter trees with non-uniform weights (average case)}
\label{tab:bounded_diameter_w}
\begin{tabular}{|c|c|c|c|c|}
\hline
$n$ & Avg Cost & Approx Ratio & Avg Time (ms) & Std Dev \\
\hline
\csvreader[head to column names]{data/bounded_diameter_w.csv}{}{%
\n & \avgcost & \approxratio & \avgtime & \stddev \\
}
\hline
\end{tabular}
\end{table}
\begin{figure}[H]
\centering
\begin{tikzpicture}
\begin{axis}[
    xlabel={Tree size ($n$)},
    ylabel={Average cost},
    width=0.48\textwidth,
    height=6cm,
    grid=major,
    legend pos=north west,
    mark size=2pt
]
\addplot[blue, mark=*, thick] table[x=n, y=avgcost, col sep=comma] {data/bounded_diameter_w.csv};
\legend{Avg Cost}
\end{axis}
\end{tikzpicture}
\hfill
\begin{tikzpicture}
\begin{axis}[
    xlabel={Tree size ($n$)},
    ylabel={Runtime (ms)},
    width=0.48\textwidth,
    height=6cm,
    grid=major,
    legend pos=north west,
    mark size=2pt
]
\addplot[red, mark=square*, thick] table[x=n, y=avgtime, col sep=comma] {data/bounded_diameter_w.csv};
\legend{Avg Time}
\end{axis}
\end{tikzpicture}
\caption{Performance metrics: Bounded diameter trees with non-uniform weights}
\label{fig:bounded_diameter_w}
\end{figure}

\input{chapters/experimental_results/trees_average_case_non_uniform_weights/large_degree_w.tex}
\input{chapters/experimental_results/trees_average_case_non_uniform_weights/large_diameter_w.tex}
\subsection{Star Trees with Non-Uniform Weights ($T_* || V, w || \sum w_i C_i$)}
\label{sec:star_w}

Results for star trees with non-uniform vertex weights.

\begin{table}[H]
\centering
\caption{Star trees with non-uniform weights (average case)}
\label{tab:star_w}
\begin{tabular}{|c|c|c|c|c|}
\hline
$n$ & Avg Cost & Approx Ratio & Avg Time (ms) & Std Dev \\
\hline
\csvreader[head to column names]{data/star_w.csv}{}{%
\n & \avgcost & \approxratio & \avgtime & \stddev \\
}
\hline
\end{tabular}
\end{table}
\begin{figure}[H]
\centering
\begin{tikzpicture}
\begin{axis}[
    xlabel={Tree size ($n$)},
    ylabel={Average cost},
    width=0.48\textwidth,
    height=6cm,
    grid=major,
    legend pos=north west,
    mark size=2pt
]
\addplot[blue, mark=*, thick] table[x=n, y=avgcost, col sep=comma] {data/star_w.csv};
\legend{Avg Cost}
\end{axis}
\end{tikzpicture}
\hfill
\begin{tikzpicture}
\begin{axis}[
    xlabel={Tree size ($n$)},
    ylabel={Runtime (ms)},
    width=0.48\textwidth,
    height=6cm,
    grid=major,
    legend pos=north west,
    mark size=2pt
]
\addplot[red, mark=square*, thick] table[x=n, y=avgtime, col sep=comma] {data/star_w.csv};
\legend{Avg Time}
\end{axis}
\end{tikzpicture}
\caption{Performance metrics: Star trees with non-uniform weights}
\label{fig:star_w}
\end{figure}

\subsection{Spider Trees with Non-Uniform Weights ($T_{sp} || V, w || \sum w_i C_i$)}
\label{sec:spider_w}

Results for spider trees with non-uniform vertex weights.

\begin{table}[H]
\centering
\caption{Spider trees with non-uniform weights (average case)}
\label{tab:spider_w}
\begin{tabular}{|c|c|c|c|c|}
\hline
$n$ & Avg Cost & Approx Ratio & Avg Time (ms) & Std Dev \\
\hline
\csvreader[head to column names]{data/spider_w.csv}{}{%
\n & \avgcost & \approxratio & \avgtime & \stddev \\
}
\hline
\end{tabular}
\end{table}
\begin{figure}[H]
\centering
\begin{tikzpicture}
\begin{axis}[
    xlabel={Tree size ($n$)},
    ylabel={Average cost},
    width=0.48\textwidth,
    height=6cm,
    grid=major,
    legend pos=north west,
    mark size=2pt
]
\addplot[blue, mark=*, thick] table[x=n, y=avgcost, col sep=comma] {data/spider_w.csv};
\legend{Avg Cost}
\end{axis}
\end{tikzpicture}
\hfill
\begin{tikzpicture}
\begin{axis}[
    xlabel={Tree size ($n$)},
    ylabel={Runtime (ms)},
    width=0.48\textwidth,
    height=6cm,
    grid=major,
    legend pos=north west,
    mark size=2pt
]
\addplot[red, mark=square*, thick] table[x=n, y=avgtime, col sep=comma] {data/spider_w.csv};
\legend{Avg Time}
\end{axis}
\end{tikzpicture}
\caption{Performance metrics: Spider trees with non-uniform weights}
\label{fig:spider_w}
\end{figure}

\input{chapters/experimental_results/trees_average_case_non_uniform_weights/caterpillar_w.tex}
\subsection{Balanced Trees with Non-Uniform Weights ($T_{bal} || V, w || \sum w_i C_i$)}
\label{sec:balanced_w}

Results for balanced trees with non-uniform vertex weights.

\begin{table}[H]
\centering
\caption{Balanced trees with non-uniform weights (average case)}
\label{tab:balanced_w}
\begin{tabular}{|c|c|c|c|c|}
\hline
$n$ & Avg Cost & Approx Ratio & Avg Time (ms) & Std Dev \\
\hline
\csvreader[head to column names]{data/balanced_w.csv}{}{%
\n & \avgcost & \approxratio & \avgtime & \stddev \\
}
\hline
\end{tabular}
\end{table}
\begin{figure}[H]
\centering
\begin{tikzpicture}
\begin{axis}[
    xlabel={Tree size ($n$)},
    ylabel={Average cost},
    width=0.48\textwidth,
    height=6cm,
    grid=major,
    legend pos=north west,
    mark size=2pt
]
\addplot[blue, mark=*, thick] table[x=n, y=avgcost, col sep=comma] {data/balanced_w.csv};
\legend{Avg Cost}
\end{axis}
\end{tikzpicture}
\hfill
\begin{tikzpicture}
\begin{axis}[
    xlabel={Tree size ($n$)},
    ylabel={Runtime (ms)},
    width=0.48\textwidth,
    height=6cm,
    grid=major,
    legend pos=north west,
    mark size=2pt
]
\addplot[red, mark=square*, thick] table[x=n, y=avgtime, col sep=comma] {data/balanced_w.csv};
\legend{Avg Time}
\end{axis}
\end{tikzpicture}
\caption{Performance metrics: Balanced trees with non-uniform weights}
\label{fig:balanced_w}
\end{figure}

\subsection{Broom Trees with Non-Uniform Weights ($T_{br} || V, w || \sum w_i C_i$)}
\label{sec:broom_w}

Results for broom trees with non-uniform vertex weights.

\begin{table}[H]
\centering
\caption{Broom trees with non-uniform weights (average case)}
\label{tab:broom_w}
\begin{tabular}{|c|c|c|c|c|}
\hline
$n$ & Avg Cost & Approx Ratio & Avg Time (ms) & Std Dev \\
\hline
\csvreader[head to column names]{data/broom_w.csv}{}{%
\n & \avgcost & \approxratio & \avgtime & \stddev \\
}
\hline
\end{tabular}
\end{table}
\begin{figure}[H]
\centering
\begin{tikzpicture}
\begin{axis}[
    xlabel={Tree size ($n$)},
    ylabel={Average cost},
    width=0.48\textwidth,
    height=6cm,
    grid=major,
    legend pos=north west,
    mark size=2pt
]
\addplot[blue, mark=*, thick] table[x=n, y=avgcost, col sep=comma] {data/broom_w.csv};
\legend{Avg Cost}
\end{axis}
\end{tikzpicture}
\hfill
\begin{tikzpicture}
\begin{axis}[
    xlabel={Tree size ($n$)},
    ylabel={Runtime (ms)},
    width=0.48\textwidth,
    height=6cm,
    grid=major,
    legend pos=north west,
    mark size=2pt
]
\addplot[red, mark=square*, thick] table[x=n, y=avgtime, col sep=comma] {data/broom_w.csv};
\legend{Avg Time}
\end{axis}
\end{tikzpicture}
\caption{Performance metrics: Broom trees with non-uniform weights}
\label{fig:broom_w}
\end{figure}

\subsection{Lobster Trees with Non-Uniform Weights ($T_{lob} || V, w || \sum w_i C_i$)}
\label{sec:lobster_w}

Results for lobster trees with non-uniform vertex weights.

\begin{table}[H]
\centering
\caption{Lobster trees with non-uniform weights (average case)}
\label{tab:lobster_w}
\begin{tabular}{|c|c|c|c|c|}
\hline
$n$ & Avg Cost & Approx Ratio & Avg Time (ms) & Std Dev \\
\hline
\csvreader[head to column names]{data/lobster_w.csv}{}{%
\n & \avgcost & \approxratio & \avgtime & \stddev \\
}
\hline
\end{tabular}
\end{table}
\begin{figure}[H]
\centering
\begin{tikzpicture}
\begin{axis}[
    xlabel={Tree size ($n$)},
    ylabel={Average cost},
    width=0.48\textwidth,
    height=6cm,
    grid=major,
    legend pos=north west,
    mark size=2pt
]
\addplot[blue, mark=*, thick] table[x=n, y=avgcost, col sep=comma] {data/lobster_w.csv};
\legend{Avg Cost}
\end{axis}
\end{tikzpicture}
\hfill
\begin{tikzpicture}
\begin{axis}[
    xlabel={Tree size ($n$)},
    ylabel={Runtime (ms)},
    width=0.48\textwidth,
    height=6cm,
    grid=major,
    legend pos=north west,
    mark size=2pt
]
\addplot[red, mark=square*, thick] table[x=n, y=avgtime, col sep=comma] {data/lobster_w.csv};
\legend{Avg Time}
\end{axis}
\end{tikzpicture}
\caption{Performance metrics: Lobster trees with non-uniform weights}
\label{fig:lobster_w}
\end{figure}

\subsection{Full m-ary Trees with Non-Uniform Weights ($T_m || V, w || \sum w_i C_i$)}
\label{sec:mary_w}

Results for full m-ary trees with non-uniform vertex weights.

\begin{table}[H]
\centering
\caption{Full m-ary trees with non-uniform weights (average case)}
\label{tab:mary_w}
\begin{tabular}{|c|c|c|c|c|}
\hline
$n$ & Avg Cost & Approx Ratio & Avg Time (ms) & Std Dev \\
\hline
\csvreader[head to column names]{data/mary_w.csv}{}{%
\n & \avgcost & \approxratio & \avgtime & \stddev \\
}
\hline
\end{tabular}
\end{table}
\begin{figure}[H]
\centering
\begin{tikzpicture}
\begin{axis}[
    xlabel={Tree size ($n$)},
    ylabel={Average cost},
    width=0.48\textwidth,
    height=6cm,
    grid=major,
    legend pos=north west,
    mark size=2pt
]
\addplot[blue, mark=*, thick] table[x=n, y=avgcost, col sep=comma] {data/mary_w.csv};
\legend{Avg Cost}
\end{axis}
\end{tikzpicture}
\hfill
\begin{tikzpicture}
\begin{axis}[
    xlabel={Tree size ($n$)},
    ylabel={Runtime (ms)},
    width=0.48\textwidth,
    height=6cm,
    grid=major,
    legend pos=north west,
    mark size=2pt
]
\addplot[red, mark=square*, thick] table[x=n, y=avgtime, col sep=comma] {data/mary_w.csv};
\legend{Avg Time}
\end{axis}
\end{tikzpicture}
\caption{Performance metrics: Full m-ary trees with non-uniform weights}
\label{fig:mary_w}
\end{figure}

\input{chapters/experimental_results/trees_average_case_non_uniform_weights/probabilistic_w.tex}
\input{chapters/experimental_results/trees_average_case_non_uniform_weights/recursive_w.tex}
\subsection{Preferential Attachment Trees with Non-Uniform Weights ($T_{PA} || V, w || \sum w_i C_i$)}
\label{sec:preferential_w}

Results for preferential attachment trees with non-uniform vertex weights.

\begin{table}[H]
\centering
\caption{Preferential attachment trees with non-uniform weights (average case)}
\label{tab:preferential_w}
\begin{tabular}{|c|c|c|c|c|}
\hline
$n$ & Avg Cost & Approx Ratio & Avg Time (ms) & Std Dev \\
\hline
\csvreader[head to column names]{data/preferential_w.csv}{}{%
\n & \avgcost & \approxratio & \avgtime & \stddev \\
}
\hline
\end{tabular}
\end{table}
\begin{figure}[H]
\centering
\begin{tikzpicture}
\begin{axis}[
    xlabel={Tree size ($n$)},
    ylabel={Average cost},
    width=0.48\textwidth,
    height=6cm,
    grid=major,
    legend pos=north west,
    mark size=2pt
]
\addplot[blue, mark=*, thick] table[x=n, y=avgcost, col sep=comma] {data/preferential_w.csv};
\legend{Avg Cost}
\end{axis}
\end{tikzpicture}
\hfill
\begin{tikzpicture}
\begin{axis}[
    xlabel={Tree size ($n$)},
    ylabel={Runtime (ms)},
    width=0.48\textwidth,
    height=6cm,
    grid=major,
    legend pos=north west,
    mark size=2pt
]
\addplot[red, mark=square*, thick] table[x=n, y=avgtime, col sep=comma] {data/preferential_w.csv};
\legend{Avg Time}
\end{axis}
\end{tikzpicture}
\caption{Performance metrics: Preferential attachment trees with non-uniform weights}
\label{fig:preferential_w}
\end{figure}



\clearpage
\section{Trees with Non-Uniform Costs (Worst Case)}

This section evaluates algorithms for the worst-case search problem with non-uniform query costs. Due to the limited computational resources and the increased complexity of handling non-uniform costs, for each algorithm, we include results up to the largest tree size that could be processed within the available time. Since we are interested in the approximation ratios achieves by the algorithms, for small instances for which we were able to compute the optimal solution, we include comparisons against the optimum as well. If such results are not available for a particular tree size, it is indicated in the respective tables by the N/A entry.

As the running times of the algorithms vary significantly, all plots of the running times use a logarithmic scale on the y-axis to better visualize these differences. Also, since the running times are often polynomial os high degree it would be hard to distinguish their running times from the exponential time procedures on the linear scale.

\subsection{Random Trees with Uniform Cost Distribution ($T_R || V, c \sim \mathcal{U} || \COST$)}
\label{sec:random_c_uniform}

\subsubsection{Exact Algorithm}

\begin{table}[ht]
\centering
\tiny
\caption{Random trees with uniform cost distribution - exact algorithm}
\label{tab:random_c_uniform_dp}
\begin{tabular}{|c|c|c|}
\hline
$n$ & Avg Time (s) & Max Time (s) \\
\hline
\csvreader[head to column names, late after line=\\]{data/dp_tree/dp_tree_basic_cu_analysis.csv}{Size=\Size,AvgCPUTime_s=\AvgTime,MaxCPUTime_s=\MaxTime,AvgCost=\AvgCost,MaxCost=\MaxCost,InstanceCount=\Instances}{%
\Size & \AvgTime & \MaxTime}
\hline
\end{tabular}
\end{table}

\begin{figure}[ht]
\centering
\begin{tikzpicture}
\begin{axis}[
    xlabel={Tree size ($n$)},
    ylabel={Runtime (s)},
    ymode=log,
    width=\textwidth,
    height=7cm,
    grid=major,
    legend pos=north west,
    mark size=2pt
]
\addplot[blue, mark=*, thick] table[x=Size, y=AvgCPUTime_s, col sep=comma] {data/dp_tree/dp_tree_basic_cu_analysis.csv};
\addplot[red, mark=square*, thick] table[x=Size, y=MaxCPUTime_s, col sep=comma] {data/dp_tree/dp_tree_basic_cu_analysis.csv};
\legend{Avg Time, Max Time}
\end{axis}
\end{tikzpicture}
\caption{Random trees with uniform costs - exact algorithm: runtime}
\label{fig:random_c_uniform_dp_time}
\end{figure}

\FloatBarrier

\subsubsection{$O(\log n / \log \log n)$-Approximation Algorithm}

\begin{table}[ht]
\centering
\tiny
\caption{Random trees with uniform cost distribution - $O(\log n / \log \log n)$-approximation algorithm}
\label{tab:random_c_uniform_cic}
\begin{tabular}{|c|c|c|c|c|}
\hline
$n$ & Avg Time (s) & Max Time (s) & Avg Ratio & Max Ratio \\
\hline
\csvreader[head to column names, late after line=\\]{data/cic_insp/cic_insp_basic_cu_comparison.csv}{Size=\Size,AvgCPUTime_s=\AvgTime,MaxCPUTime_s=\MaxTime,AvgCost=\AvgCost,MaxCost=\MaxCost,AvgRatio=\AvgRatio,MaxRatio=\MaxRatio}{%
\Size & \AvgTime & \MaxTime & \AvgRatio & \MaxRatio}
\hline
\end{tabular}
\end{table}

\begin{figure}[ht]
\centering
\begin{tikzpicture}
\begin{axis}[
    xlabel={Tree size ($n$)},
    ylabel={Approximation ratio},
    width=\textwidth,
    height=7cm,
    grid=major,
    legend pos=north west,
    mark size=2pt
]
\addplot[blue, mark=*, thick] table[x=Size, y=AvgRatio, col sep=comma] {data/cic_insp/cic_insp_basic_cu_comparison.csv};
\addplot[red, mark=square*, thick] table[x=Size, y=MaxRatio, col sep=comma] {data/cic_insp/cic_insp_basic_cu_comparison.csv};
\legend{Avg Ratio, Max Ratio}
\end{axis}
\end{tikzpicture}
\caption{Random trees with uniform costs - $O(\log n / \log \log n)$-approximation algorithm: approximation quality}
\label{fig:random_c_uniform_cic}
\end{figure}

\begin{figure}[ht]
\centering
\begin{tikzpicture}
\begin{axis}[
    xlabel={Tree size ($n$)},
    ylabel={Runtime (s)},
    ymode=log,
    width=\textwidth,
    height=7cm,
    grid=major,
    legend pos=north west,
    mark size=2pt
]
\addplot[blue, mark=*, thick] table[x=Size, y=AvgCPUTime_s, col sep=comma] {data/cic_insp/cic_insp_basic_cu_comparison.csv};
\addplot[red, mark=square*, thick] table[x=Size, y=MaxCPUTime_s, col sep=comma] {data/cic_insp/cic_insp_basic_cu_comparison.csv};
\legend{Avg Time, Max Time}
\end{axis}
\end{tikzpicture}
\caption{Random trees with uniform costs - $O(\log n / \log \log n)$-approximation algorithm: runtime}
\label{fig:random_c_uniform_cic_time}
\end{figure}

\FloatBarrier

\subsubsection{$O(\sqrt{\log n})$-Approximation Algorithm}

\begin{table}[ht]
\centering
\tiny
\caption{Random trees with uniform cost distribution - $O(\sqrt{\log n})$-approximation algorithm}
\label{tab:random_c_uniform_der}
\begin{tabular}{|c|c|c|c|c|}
\hline
$n$ & Avg Time (s) & Max Time (s) & Avg Ratio & Max Ratio \\
\hline
\csvreader[head to column names, late after line=\\]{data/der_insp/der_insp_basic_cu_comparison.csv}{Size=\Size,AvgCPUTime_s=\AvgTime,MaxCPUTime_s=\MaxTime,AvgCost=\AvgCost,MaxCost=\MaxCost,AvgRatio=\AvgRatio,MaxRatio=\MaxRatio}{%
\Size & \AvgTime & \MaxTime & \AvgRatio & \MaxRatio}
\hline
\end{tabular}
\end{table}

\begin{figure}[ht]
\centering
\begin{tikzpicture}
\begin{axis}[
    xlabel={Tree size ($n$)},
    ylabel={Approximation ratio},
    width=\textwidth,
    height=7cm,
    grid=major,
    legend pos=north west,
    mark size=2pt
]
\addplot[blue, mark=*, thick] table[x=Size, y=AvgRatio, col sep=comma] {data/der_insp/der_insp_basic_cu_comparison.csv};
\addplot[red, mark=square*, thick] table[x=Size, y=MaxRatio, col sep=comma] {data/der_insp/der_insp_basic_cu_comparison.csv};
\legend{Avg Ratio, Max Ratio}
\end{axis}
\end{tikzpicture}
\caption{Random trees with uniform costs - $O(\sqrt{\log n})$-approximation algorithm: approximation quality}
\label{fig:random_c_uniform_der}
\end{figure}

\begin{figure}[ht]
\centering
\begin{tikzpicture}
\begin{axis}[
    xlabel={Tree size ($n$)},
    ylabel={Runtime (s)},
    ymode=log,
    width=\textwidth,
    height=7cm,
    grid=major,
    legend pos=north west,
    mark size=2pt
]
\addplot[blue, mark=*, thick] table[x=Size, y=AvgCPUTime_s, col sep=comma] {data/der_insp/der_insp_basic_cu_comparison.csv};
\addplot[red, mark=square*, thick] table[x=Size, y=MaxCPUTime_s, col sep=comma] {data/der_insp/der_insp_basic_cu_comparison.csv};
\legend{Avg Time, Max Time}
\end{axis}
\end{tikzpicture}
\caption{Random trees with uniform costs - $O(\sqrt{\log n})$-approximation algorithm: runtime}
\label{fig:random_c_uniform_der_time}
\end{figure}

\FloatBarrier

\subsubsection{$O(\log \log n)$-Approximation Algorithm}

\begin{table}[ht]
\centering
\tiny
\caption{Random trees with uniform cost distribution - $O(\log \log n)$-approximation algorithm}
\label{tab:random_c_uniform_kup}
\begin{tabular}{|c|c|c|c|c|}
\hline
$n$ & Avg Time (s) & Max Time (s) & Avg Ratio & Max Ratio \\
\hline
\csvreader[head to column names, late after line=\\]{data/k_up/k_up_basic_cu_comparison.csv}{Size=\Size,AvgCPUTime_s=\AvgTime,MaxCPUTime_s=\MaxTime,AvgCost=\AvgCost,MaxCost=\MaxCost,AvgRatio=\AvgRatio,MaxRatio=\MaxRatio}{%
\Size & \AvgTime & \MaxTime & \AvgRatio & \MaxRatio}
\hline
\end{tabular}
\end{table}

\begin{figure}[ht]
\centering
\begin{tikzpicture}
\begin{axis}[
    xlabel={Tree size ($n$)},
    ylabel={Approximation ratio},
    width=\textwidth,
    height=7cm,
    grid=major,
    legend pos=north west,
    mark size=2pt
]
\addplot[blue, mark=*, thick] table[x=Size, y=AvgRatio, col sep=comma] {data/k_up/k_up_basic_cu_comparison.csv};
\addplot[red, mark=square*, thick] table[x=Size, y=MaxRatio, col sep=comma] {data/k_up/k_up_basic_cu_comparison.csv};
\legend{Avg Ratio, Max Ratio}
\end{axis}
\end{tikzpicture}
\caption{Random trees with uniform costs - $O(\log \log n)$-approximation algorithm: approximation quality}
\label{fig:random_c_uniform_kup}
\end{figure}

\begin{figure}[ht]
\centering
\begin{tikzpicture}
\begin{axis}[
    xlabel={Tree size ($n$)},
    ylabel={Runtime (s)},
    ymode=log,
    width=\textwidth,
    height=7cm,
    grid=major,
    legend pos=north west,
    mark size=2pt
]
\addplot[blue, mark=*, thick] table[x=Size, y=AvgCPUTime_s, col sep=comma] {data/k_up/k_up_basic_cu_comparison.csv};
\addplot[red, mark=square*, thick] table[x=Size, y=MaxCPUTime_s, col sep=comma] {data/k_up/k_up_basic_cu_comparison.csv};
\legend{Avg Time, Max Time}
\end{axis}
\end{tikzpicture}
\caption{Random trees with uniform costs - $O(\log \log n)$-approximation algorithm: runtime}
\label{fig:random_c_uniform_kup_time}
\end{figure}

\FloatBarrier
We have the following observations:
\begin{itemize}
    \item We see that the running time of the exact procedure plotted on the logarithmic scale behaves as a linear function. This matches our expectation since it means that the running time is exponential in $n$. Notably, this running time is lower then the worst case running time of $O\br{2^nn}$, suggesting that for our instances the base of the exponential function representing the running time of the algorithm is smaller than 2.
    \item Among the approximation algorothm the best running time was always achieved by the $O\br{\log n/\log \log n}$-approximation algorithm. 
    \item The approximation algorithms show a polynomial growth in running time, with the $O(\log n / \log \log n)$-approximation algorithm being the fastest among them, followed by the $O(\sqrt{\log n})$-approximation algorithm, and finally the $O(\log \log n)$-approximation algorithm being the slowest.
    \item The approximation ratios achieved by all approximation algorithms are significantly better than their theoretical worst-case guarantees. All of the approximation ratios were below 2. We also see that no algorithm consistently outperformed the other.
    \item  All of the above observations will also be prevalent among further experiments except when noted.
\end{itemize}


\clearpage\subsection{Random Trees with Binomial Cost Distribution ($T_R || V, c \sim \mathcal{B} || \COST_{max}$)}
\label{sec:random_c_binomial}

\subsubsection{Exact Algorithm}

\begin{table}[ht]
\centering
\tiny
\caption{Random trees with binomial cost distribution - exact algorithm}
\label{tab:random_c_binomial_dp}
\begin{tabular}{|c|c|c|}
\hline
$n$ & Avg Time (s) & Max Time (s) \\
\hline
\csvreader[head to column names, late after line=\\]{data/dp_tree/dp_tree_basic_cb_analysis.csv}{Size=\Size,AvgCPUTime_s=\AvgTime,MaxCPUTime_s=\MaxTime,AvgCost=\AvgCost,MaxCost=\MaxCost,InstanceCount=\Instances}{%
\Size & \AvgTime & \MaxTime}
\hline
\end{tabular}
\end{table}

\begin{figure}[ht]
\centering
\begin{tikzpicture}
\begin{axis}[
    xlabel={Tree size ($n$)},
    ylabel={Runtime (s)},
    ymode=log,
    width=\textwidth,
    height=7cm,
    grid=major,
    legend pos=north west,
    mark size=2pt
]
\addplot[blue, mark=*, thick] table[x=Size, y=AvgCPUTime_s, col sep=comma] {data/dp_tree/dp_tree_basic_cb_analysis.csv};
\addplot[red, mark=square*, thick] table[x=Size, y=MaxCPUTime_s, col sep=comma] {data/dp_tree/dp_tree_basic_cb_analysis.csv};
\legend{Avg Time, Max Time}
\end{axis}
\end{tikzpicture}
\caption{Random trees with binomial costs - exact algorithm: runtime comparison}
\label{fig:random_c_binomial_dp_time}
\end{figure}

\FloatBarrier

\subsubsection{$O(\log n / \log \log n)$-Approximation Algorithm}

\begin{table}[ht]
\centering
\tiny
\caption{Random trees with binomial cost distribution - $O(\log n / \log \log n)$-approximation algorithm}
\label{tab:random_c_binomial_cic}
\begin{tabular}{|c|c|c|c|c|}
\hline
$n$ & Avg Time (s) & Max Time (s) & Avg Ratio & Max Ratio \\
\hline
\csvreader[head to column names, late after line=\\]{data/cic_insp/cic_insp_basic_cb_comparison.csv}{Size=\Size,AvgCPUTime_s=\AvgTime,MaxCPUTime_s=\MaxTime,AvgCost=\AvgCost,MaxCost=\MaxCost,AvgRatio=\AvgRatio,MaxRatio=\MaxRatio}{%
\Size & \AvgTime & \MaxTime & \AvgRatio & \MaxRatio}
\hline
\end{tabular}
\end{table}

\begin{figure}[ht]
\centering
\begin{tikzpicture}
\begin{axis}[
    xlabel={Tree size ($n$)},
    ylabel={Runtime (s)},
    ymode=log,
    width=\textwidth,
    height=7cm,
    grid=major,
    legend pos=north west,
    mark size=2pt
]
\addplot[blue, mark=*, thick] table[x=Size, y=AvgCPUTime_s, col sep=comma] {data/cic_insp/cic_insp_basic_cb_comparison.csv};
\addplot[red, mark=square*, thick] table[x=Size, y=MaxCPUTime_s, col sep=comma] {data/cic_insp/cic_insp_basic_cb_comparison.csv};
\legend{Avg Time, Max Time}
\end{axis}
\end{tikzpicture}
\caption{Random trees with binomial costs - $O(\log n / \log \log n)$-approximation algorithm: runtime}
\label{fig:random_c_binomial_cic_time}
\end{figure}

\begin{figure}[ht]
\centering
\begin{tikzpicture}
\begin{axis}[
    xlabel={Tree size ($n$)},
    ylabel={Approximation ratio},
    width=\textwidth,
    height=7cm,
    grid=major,
    legend pos=north west,
    mark size=2pt
]
\addplot[blue, mark=*, thick] table[x=Size, y=AvgRatio, col sep=comma] {data/cic_insp/cic_insp_basic_cb_comparison.csv};
\addplot[red, mark=square*, thick] table[x=Size, y=MaxRatio, col sep=comma] {data/cic_insp/cic_insp_basic_cb_comparison.csv};
\legend{Avg Ratio, Max Ratio}
\end{axis}
\end{tikzpicture}
\caption{Random trees with binomial costs - $O(\log n / \log \log n)$-approximation algorithm: approximation quality}
\label{fig:random_c_binomial_cic}
\end{figure}

\FloatBarrier

\subsubsection{$O(\sqrt{\log n})$-Approximation Algorithm}

\begin{table}[ht]
\centering
\tiny
\caption{Random trees with binomial cost distribution - $O(\sqrt{\log n})$-approximation algorithm}
\label{tab:random_c_binomial_der}
\begin{tabular}{|c|c|c|c|c|}
\hline
$n$ & Avg Time (s) & Max Time (s) & Avg Ratio & Max Ratio \\
\hline
\csvreader[head to column names, late after line=\\]{data/der_insp/der_insp_basic_cb_comparison.csv}{Size=\Size,AvgCPUTime_s=\AvgTime,MaxCPUTime_s=\MaxTime,AvgCost=\AvgCost,MaxCost=\MaxCost,AvgRatio=\AvgRatio,MaxRatio=\MaxRatio}{%
\Size & \AvgTime & \MaxTime & \AvgRatio & \MaxRatio}
\hline
\end{tabular}
\end{table}

\begin{figure}[ht]
\centering
\begin{tikzpicture}
\begin{axis}[
    xlabel={Tree size ($n$)},
    ylabel={Runtime (s)},
    ymode=log,
    width=\textwidth,
    height=7cm,
    grid=major,
    legend pos=north west,
    mark size=2pt
]
\addplot[blue, mark=*, thick] table[x=Size, y=AvgCPUTime_s, col sep=comma] {data/der_insp/der_insp_basic_cb_comparison.csv};
\addplot[red, mark=square*, thick] table[x=Size, y=MaxCPUTime_s, col sep=comma] {data/der_insp/der_insp_basic_cb_comparison.csv};
\legend{Avg Time, Max Time}
\end{axis}
\end{tikzpicture}
\caption{Random trees with binomial costs - $O(\sqrt{\log n})$-approximation algorithm: runtime}
\label{fig:random_c_binomial_der_time}
\end{figure}

\begin{figure}[ht]
\centering
\begin{tikzpicture}
\begin{axis}[
    xlabel={Tree size ($n$)},
    ylabel={Approximation ratio},
    width=\textwidth,
    height=7cm,
    grid=major,
    legend pos=north west,
    mark size=2pt
]
\addplot[blue, mark=*, thick] table[x=Size, y=AvgRatio, col sep=comma] {data/der_insp/der_insp_basic_cb_comparison.csv};
\addplot[red, mark=square*, thick] table[x=Size, y=MaxRatio, col sep=comma] {data/der_insp/der_insp_basic_cb_comparison.csv};
\legend{Avg Ratio, Max Ratio}
\end{axis}
\end{tikzpicture}
\caption{Random trees with binomial costs - $O(\sqrt{\log n})$-approximation algorithm: approximation quality}
\label{fig:random_c_binomial_der}
\end{figure}

\FloatBarrier

\subsubsection{$O(\log \log n)$-Approximation Algorithm}

\begin{table}[ht]
\centering
\tiny
\caption{Random trees with binomial cost distribution - $O(\log \log n)$-approximation algorithm}
\label{tab:random_c_binomial_kup}
\begin{tabular}{|c|c|c|c|c|}
\hline
$n$ & Avg Time (s) & Max Time (s) & Avg Ratio & Max Ratio \\
\hline
\csvreader[head to column names, late after line=\\]{data/k_up/k_up_basic_cb_comparison.csv}{Size=\Size,AvgCPUTime_s=\AvgTime,MaxCPUTime_s=\MaxTime,AvgCost=\AvgCost,MaxCost=\MaxCost,AvgRatio=\AvgRatio,MaxRatio=\MaxRatio}{%
\Size & \AvgTime & \MaxTime & \AvgRatio & \MaxRatio}
\hline
\end{tabular}
\end{table}

\begin{figure}[ht]
\centering
\begin{tikzpicture}
\begin{axis}[
    xlabel={Tree size ($n$)},
    ylabel={Runtime (s)},
    ymode=log,
    width=\textwidth,
    height=7cm,
    grid=major,
    legend pos=north west,
    mark size=2pt
]
\addplot[blue, mark=*, thick] table[x=Size, y=AvgCPUTime_s, col sep=comma] {data/k_up/k_up_basic_cb_comparison.csv};
\addplot[red, mark=square*, thick] table[x=Size, y=MaxCPUTime_s, col sep=comma] {data/k_up/k_up_basic_cb_comparison.csv};
\legend{Avg Time, Max Time}
\end{axis}
\end{tikzpicture}
\caption{Random trees with binomial costs - $O(\log \log n)$-approximation algorithm: runtime}
\label{fig:random_c_binomial_kup_time}
\end{figure}

\begin{figure}[ht]
\centering
\begin{tikzpicture}
\begin{axis}[
    xlabel={Tree size ($n$)},
    ylabel={Approximation ratio},
    width=\textwidth,
    height=7cm,
    grid=major,
    legend pos=north west,
    mark size=2pt
]
\addplot[blue, mark=*, thick] table[x=Size, y=AvgRatio, col sep=comma] {data/k_up/k_up_basic_cb_comparison.csv};
\addplot[red, mark=square*, thick] table[x=Size, y=MaxRatio, col sep=comma] {data/k_up/k_up_basic_cb_comparison.csv};
\legend{Avg Ratio, Max Ratio}
\end{axis}
\end{tikzpicture}
\caption{Random trees with binomial costs - $O(\log \log n)$-approximation algorithm: approximation quality}
\label{fig:random_c_binomial_kup}
\end{figure}


\clearpage\subsection{Random Trees with Exponential Cost Distribution ($T_R || V, c \sim \mathcal{E} || \COST_{max}$)}
\label{sec:random_c_exponential}

\subsubsection{Exact Algorithm}

\begin{table}[ht]
\centering
\tiny
\caption{Random trees with exponential cost distribution - exact algorithm}
\label{tab:random_c_exponential_dp}
\begin{tabular}{|c|c|c|}
\hline
$n$ & Avg Time (s) & Max Time (s) \\
\hline
\csvreader[head to column names, late after line=\\]{data/dp_tree/dp_tree_basic_ce_analysis.csv}{Size=\Size,AvgCPUTime_s=\AvgTime,MaxCPUTime_s=\MaxTime,AvgCost=\AvgCost,MaxCost=\MaxCost,InstanceCount=\Instances}{%
\Size & \AvgTime & \MaxTime}
\hline
\end{tabular}
\end{table}

\begin{figure}[ht]
\centering
\begin{tikzpicture}
\begin{axis}[
    xlabel={Tree size ($n$)},
    ylabel={Runtime (s)},
    ymode=log,
    width=\textwidth,
    height=7cm,
    grid=major,
    legend pos=north west,
    mark size=2pt
]
\addplot[blue, mark=*, thick] table[x=Size, y=AvgCPUTime_s, col sep=comma] {data/dp_tree/dp_tree_basic_ce_analysis.csv};
\addplot[red, mark=square*, thick] table[x=Size, y=MaxCPUTime_s, col sep=comma] {data/dp_tree/dp_tree_basic_ce_analysis.csv};
\legend{Avg Time, Max Time}
\end{axis}
\end{tikzpicture}
\caption{Random trees with exponential costs - exact algorithm: runtime comparison}
\label{fig:random_c_exponential_dp_time}
\end{figure}

\FloatBarrier

\subsubsection{$O(\log n / \log \log n)$-Approximation Algorithm}

\begin{table}[ht]
\centering
\tiny
\caption{Random trees with exponential cost distribution - $O(\log n / \log \log n)$-approximation algorithm}
\label{tab:random_c_exponential_cic}
\begin{tabular}{|c|c|c|c|c|}
\hline
$n$ & Avg Time (s) & Max Time (s) & Avg Ratio & Max Ratio \\
\hline
\csvreader[head to column names, late after line=\\]{data/cic_insp/cic_insp_basic_ce_comparison.csv}{Size=\Size,AvgCPUTime_s=\AvgTime,MaxCPUTime_s=\MaxTime,AvgCost=\AvgCost,MaxCost=\MaxCost,AvgRatio=\AvgRatio,MaxRatio=\MaxRatio}{%
\Size & \AvgTime & \MaxTime & \AvgRatio & \MaxRatio}
\hline
\end{tabular}
\end{table}

\begin{figure}[ht]
\centering
\begin{tikzpicture}
\begin{axis}[
    xlabel={Tree size ($n$)},
    ylabel={Runtime (s)},
    ymode=log,
    width=\textwidth,
    height=7cm,
    grid=major,
    legend pos=north west,
    mark size=2pt
]
\addplot[blue, mark=*, thick] table[x=Size, y=AvgCPUTime_s, col sep=comma] {data/cic_insp/cic_insp_basic_ce_comparison.csv};
\addplot[red, mark=square*, thick] table[x=Size, y=MaxCPUTime_s, col sep=comma] {data/cic_insp/cic_insp_basic_ce_comparison.csv};
\legend{Avg Time, Max Time}
\end{axis}
\end{tikzpicture}
\caption{Random trees with exponential costs - $O(\log n / \log \log n)$-approximation algorithm: runtime}
\label{fig:random_c_exponential_cic_time}
\end{figure}

\begin{figure}[ht]
\centering
\begin{tikzpicture}
\begin{axis}[
    xlabel={Tree size ($n$)},
    ylabel={Approximation ratio},
    width=\textwidth,
    height=7cm,
    grid=major,
    legend pos=north west,
    mark size=2pt
]
\addplot[blue, mark=*, thick] table[x=Size, y=AvgRatio, col sep=comma] {data/cic_insp/cic_insp_basic_ce_comparison.csv};
\addplot[red, mark=square*, thick] table[x=Size, y=MaxRatio, col sep=comma] {data/cic_insp/cic_insp_basic_ce_comparison.csv};
\legend{Avg Ratio, Max Ratio}
\end{axis}
\end{tikzpicture}
\caption{Random trees with exponential costs - $O(\log n / \log \log n)$-approximation algorithm: approximation quality}
\label{fig:random_c_exponential_cic}
\end{figure}

\FloatBarrier

\subsubsection{$O(\sqrt{\log n})$-Approximation Algorithm}

\begin{table}[ht]
\centering
\tiny
\caption{Random trees with exponential cost distribution - $O(\sqrt{\log n})$-approximation algorithm}
\label{tab:random_c_exponential_der}
\begin{tabular}{|c|c|c|c|c|}
\hline
$n$ & Avg Time (s) & Max Time (s) & Avg Ratio & Max Ratio \\
\hline
\csvreader[head to column names, late after line=\\]{data/der_insp/der_insp_basic_ce_comparison.csv}{Size=\Size,AvgCPUTime_s=\AvgTime,MaxCPUTime_s=\MaxTime,AvgCost=\AvgCost,MaxCost=\MaxCost,AvgRatio=\AvgRatio,MaxRatio=\MaxRatio}{%
\Size & \AvgTime & \MaxTime & \AvgRatio & \MaxRatio}
\hline
\end{tabular}
\end{table}

\begin{figure}[ht]
\centering
\begin{tikzpicture}
\begin{axis}[
    xlabel={Tree size ($n$)},
    ylabel={Runtime (s)},
    ymode=log,
    width=\textwidth,
    height=7cm,
    grid=major,
    legend pos=north west,
    mark size=2pt
]
\addplot[blue, mark=*, thick] table[x=Size, y=AvgCPUTime_s, col sep=comma] {data/der_insp/der_insp_basic_ce_comparison.csv};
\addplot[red, mark=square*, thick] table[x=Size, y=MaxCPUTime_s, col sep=comma] {data/der_insp/der_insp_basic_ce_comparison.csv};
\legend{Avg Time, Max Time}
\end{axis}
\end{tikzpicture}
\caption{Random trees with exponential costs - $O(\sqrt{\log n})$-approximation algorithm: runtime}
\label{fig:random_c_exponential_der_time}
\end{figure}

\begin{figure}[ht]
\centering
\begin{tikzpicture}
\begin{axis}[
    xlabel={Tree size ($n$)},
    ylabel={Approximation ratio},
    width=\textwidth,
    height=7cm,
    grid=major,
    legend pos=north west,
    mark size=2pt
]
\addplot[blue, mark=*, thick] table[x=Size, y=AvgRatio, col sep=comma] {data/der_insp/der_insp_basic_ce_comparison.csv};
\addplot[red, mark=square*, thick] table[x=Size, y=MaxRatio, col sep=comma] {data/der_insp/der_insp_basic_ce_comparison.csv};
\legend{Avg Ratio, Max Ratio}
\end{axis}
\end{tikzpicture}
\caption{Random trees with exponential costs - $O(\sqrt{\log n})$-approximation algorithm: approximation quality}
\label{fig:random_c_exponential_der}
\end{figure}

\FloatBarrier

\subsubsection{$O(\log \log n)$-Approximation Algorithm}

\begin{table}[ht]
\centering
\tiny
\caption{Random trees with exponential cost distribution - $O(\log \log n)$-approximation algorithm}
\label{tab:random_c_exponential_kup}
\begin{tabular}{|c|c|c|c|c|}
\hline
$n$ & Avg Time (s) & Max Time (s) & Avg Ratio & Max Ratio \\
\hline
\csvreader[head to column names, late after line=\\]{data/k_up/k_up_basic_ce_comparison.csv}{Size=\Size,AvgCPUTime_s=\AvgTime,MaxCPUTime_s=\MaxTime,AvgCost=\AvgCost,MaxCost=\MaxCost,AvgRatio=\AvgRatio,MaxRatio=\MaxRatio}{%
\Size & \AvgTime & \MaxTime & \AvgRatio & \MaxRatio}
\hline
\end{tabular}
\end{table}

\begin{figure}[ht]
\centering
\begin{tikzpicture}
\begin{axis}[
    xlabel={Tree size ($n$)},
    ylabel={Runtime (s)},
    ymode=log,
    width=\textwidth,
    height=7cm,
    grid=major,
    legend pos=north west,
    mark size=2pt
]
\addplot[blue, mark=*, thick] table[x=Size, y=AvgCPUTime_s, col sep=comma] {data/k_up/k_up_basic_ce_comparison.csv};
\addplot[red, mark=square*, thick] table[x=Size, y=MaxCPUTime_s, col sep=comma] {data/k_up/k_up_basic_ce_comparison.csv};
\legend{Avg Time, Max Time}
\end{axis}
\end{tikzpicture}
\caption{Random trees with exponential costs - $O(\log \log n)$-approximation algorithm: runtime}
\label{fig:random_c_exponential_kup_time}
\end{figure}

\begin{figure}[ht]
\centering
\begin{tikzpicture}
\begin{axis}[
    xlabel={Tree size ($n$)},
    ylabel={Approximation ratio},
    width=\textwidth,
    height=7cm,
    grid=major,
    legend pos=north west,
    mark size=2pt
]
\addplot[blue, mark=*, thick] table[x=Size, y=AvgRatio, col sep=comma] {data/k_up/k_up_basic_ce_comparison.csv};
\addplot[red, mark=square*, thick] table[x=Size, y=MaxRatio, col sep=comma] {data/k_up/k_up_basic_ce_comparison.csv};
\legend{Avg Ratio, Max Ratio}
\end{axis}
\end{tikzpicture}
\caption{Random trees with exponential costs - $O(\log \log n)$-approximation algorithm: approximation quality}
\label{fig:random_c_exponential_kup}
\end{figure}
\FloatBarrier
We observe that this time the $O\br{\log \log n}$-approximation algorithm encountered instances for which it produced decision tree with approximation ratio above $2.5$. Therefore we conclude that its performances was the worst among the three algorithms. The $O\br{\log n/\log \log n}$-approximation algorithm achieved the best worst approximation ratios while the $O\br{\sqrt{\log n}}$-approximation algorithm achieved the best average approximation ratios.



\clearpage\subsection{Random bounded degree trees with Uniform Cost Distribution ($T_{\Delta} || V, c \sim \mathcal{U} || \COST_{max}$)}
\label{sec:bounded_degree_c_uniform}

\subsubsection{Exact Algorithm}

\begin{table}[ht]
\centering
\tiny
\caption{Random bounded degree trees ($\Delta(T) = 3$) with uniform cost distribution - exact algorithm}
\label{tab:bounded_degree_c_uniform_dp}
\begin{tabular}{|c|c|c|}
\hline
$n$ & Avg Time (s) & Max Time (s) \\
\hline
\csvreader[head to column names, late after line=\\]{data/dp_tree/dp_tree_deg_3_cu_analysis.csv}{Size=\Size,AvgCPUTime_s=\AvgTime,MaxCPUTime_s=\MaxTime,AvgCost=\AvgCost,MaxCost=\MaxCost,InstanceCount=\Instances}{%
\Size & \AvgTime & \MaxTime}
\hline
\end{tabular}
\end{table}

\begin{figure}[ht]
\centering
\begin{tikzpicture}
\begin{axis}[
    xlabel={Tree size ($n$)},
    ylabel={Runtime (s)},
    ymode=log,
    width=\textwidth,
    height=7cm,
    grid=major,
    legend pos=north west,
    mark size=2pt
]
\addplot[blue, mark=*, thick] table[x=Size, y=AvgCPUTime_s, col sep=comma] {data/dp_tree/dp_tree_deg_3_cu_analysis.csv};
\addplot[red, mark=square*, thick] table[x=Size, y=MaxCPUTime_s, col sep=comma] {data/dp_tree/dp_tree_deg_3_cu_analysis.csv};
\legend{Avg Time, Max Time}
\end{axis}
\end{tikzpicture}
\caption{Random bounded degree trees with uniform costs - exact algorithm: runtime comparison}
\label{fig:bounded_degree_c_uniform_dp_time}
\end{figure}

\FloatBarrier

\subsubsection{$O(\log n / \log \log n)$-Approximation Algorithm}

\begin{table}[ht]
\centering
\tiny
\caption{Random bounded degree trees ($\Delta(T) = 3$) with uniform cost distribution - $O(\log n / \log \log n)$-approximation algorithm}
\label{tab:bounded_degree_c_uniform_cic}
\begin{tabular}{|c|c|c|c|c|}
\hline
$n$ & Avg Time (s) & Max Time (s) & Avg Ratio & Max Ratio \\
\hline
\csvreader[head to column names, late after line=\\]{data/cic_insp/cic_insp_deg_3_cu_comparison.csv}{Size=\Size,AvgCPUTime_s=\AvgTime,MaxCPUTime_s=\MaxTime,AvgCost=\AvgCost,MaxCost=\MaxCost,AvgRatio=\AvgRatio,MaxRatio=\MaxRatio}{%
\Size & \AvgTime & \MaxTime & \AvgRatio & \MaxRatio}
\hline
\end{tabular}
\end{table}

\begin{figure}[ht]
\centering
\begin{tikzpicture}
\begin{axis}[
    xlabel={Tree size ($n$)},
    ylabel={Runtime (s)},
    ymode=log,
    width=\textwidth,
    height=7cm,
    grid=major,
    legend pos=north west,
    mark size=2pt
]
\addplot[blue, mark=*, thick] table[x=Size, y=AvgCPUTime_s, col sep=comma] {data/cic_insp/cic_insp_deg_3_cu_comparison.csv};
\addplot[red, mark=square*, thick] table[x=Size, y=MaxCPUTime_s, col sep=comma] {data/cic_insp/cic_insp_deg_3_cu_comparison.csv};
\legend{Avg Time, Max Time}
\end{axis}
\end{tikzpicture}
\caption{Random bounded degree trees with uniform costs - $O(\log n / \log \log n)$-approximation algorithm: runtime}
\label{fig:bounded_degree_c_uniform_cic_time}
\end{figure}

\begin{figure}[ht]
\centering
\begin{tikzpicture}
\begin{axis}[
    xlabel={Tree size ($n$)},
    ylabel={Approximation ratio},
    width=\textwidth,
    height=7cm,
    grid=major,
    legend pos=north west,
    mark size=2pt
]
\addplot[blue, mark=*, thick] table[x=Size, y=AvgRatio, col sep=comma] {data/cic_insp/cic_insp_deg_3_cu_comparison.csv};
\addplot[red, mark=square*, thick] table[x=Size, y=MaxRatio, col sep=comma] {data/cic_insp/cic_insp_deg_3_cu_comparison.csv};
\legend{Avg Ratio, Max Ratio}
\end{axis}
\end{tikzpicture}
\caption{Random bounded degree trees with uniform costs - $O(\log n / \log \log n)$-approximation algorithm: approximation quality}
\label{fig:bounded_degree_c_uniform_cic}
\end{figure}

\FloatBarrier

\subsubsection{$O(\sqrt{\log n})$-Approximation Algorithm}

\begin{table}[ht]
\centering
\tiny
\caption{Random bounded degree trees ($\Delta(T) = 3$) with uniform cost distribution - $O(\sqrt{\log n})$-approximation algorithm}
\label{tab:bounded_degree_c_uniform_der}
\begin{tabular}{|c|c|c|c|c|}
\hline
$n$ & Avg Time (s) & Max Time (s) & Avg Ratio & Max Ratio \\
\hline
\csvreader[head to column names, late after line=\\]{data/der_insp/der_insp_deg_3_cu_comparison.csv}{Size=\Size,AvgCPUTime_s=\AvgTime,MaxCPUTime_s=\MaxTime,AvgCost=\AvgCost,MaxCost=\MaxCost,AvgRatio=\AvgRatio,MaxRatio=\MaxRatio}{%
\Size & \AvgTime & \MaxTime & \AvgRatio & \MaxRatio}
\hline
\end{tabular}
\end{table}

\begin{figure}[ht]
\centering
\begin{tikzpicture}
\begin{axis}[
    xlabel={Tree size ($n$)},
    ylabel={Runtime (s)},
    ymode=log,
    width=\textwidth,
    height=7cm,
    grid=major,
    legend pos=north west,
    mark size=2pt
]
\addplot[blue, mark=*, thick] table[x=Size, y=AvgCPUTime_s, col sep=comma] {data/der_insp/der_insp_deg_3_cu_comparison.csv};
\addplot[red, mark=square*, thick] table[x=Size, y=MaxCPUTime_s, col sep=comma] {data/der_insp/der_insp_deg_3_cu_comparison.csv};
\legend{Avg Time, Max Time}
\end{axis}
\end{tikzpicture}
\caption{Random bounded degree trees with uniform costs - $O(\sqrt{\log n})$-approximation algorithm: runtime}
\label{fig:bounded_degree_c_uniform_der_time}
\end{figure}

\begin{figure}[ht]
\centering
\begin{tikzpicture}
\begin{axis}[
    xlabel={Tree size ($n$)},
    ylabel={Approximation ratio},
    width=\textwidth,
    height=7cm,
    grid=major,
    legend pos=north west,
    mark size=2pt
]
\addplot[blue, mark=*, thick] table[x=Size, y=AvgRatio, col sep=comma] {data/der_insp/der_insp_deg_3_cu_comparison.csv};
\addplot[red, mark=square*, thick] table[x=Size, y=MaxRatio, col sep=comma] {data/der_insp/der_insp_deg_3_cu_comparison.csv};
\legend{Avg Ratio, Max Ratio}
\end{axis}
\end{tikzpicture}
\caption{Random bounded degree trees with uniform costs - $O(\sqrt{\log n})$-approximation algorithm: approximation quality}
\label{fig:bounded_degree_c_uniform_der}
\end{figure}

\FloatBarrier

\subsubsection{$O(\log \log n)$-Approximation Algorithm}

\begin{table}[ht]
\centering
\tiny
\caption{Random bounded degree trees ($\Delta(T) = 3$) with uniform cost distribution - $O(\log \log n)$-approximation algorithm}
\label{tab:bounded_degree_c_uniform_kup}
\begin{tabular}{|c|c|c|c|c|}
\hline
$n$ & Avg Time (s) & Max Time (s) & Avg Ratio & Max Ratio \\
\hline
\csvreader[head to column names, late after line=\\]{data/k_up/k_up_deg_3_cu_comparison.csv}{Size=\Size,AvgCPUTime_s=\AvgTime,MaxCPUTime_s=\MaxTime,AvgCost=\AvgCost,MaxCost=\MaxCost,AvgRatio=\AvgRatio,MaxRatio=\MaxRatio}{%
\Size & \AvgTime & \MaxTime & \AvgRatio & \MaxRatio}
\hline
\end{tabular}
\end{table}

\begin{figure}[ht]
\centering
\begin{tikzpicture}
\begin{axis}[
    xlabel={Tree size ($n$)},
    ylabel={Runtime (s)},
    ymode=log,
    width=\textwidth,
    height=7cm,
    grid=major,
    legend pos=north west,
    mark size=2pt
]
\addplot[blue, mark=*, thick] table[x=Size, y=AvgCPUTime_s, col sep=comma] {data/k_up/k_up_deg_3_cu_comparison.csv};
\addplot[red, mark=square*, thick] table[x=Size, y=MaxCPUTime_s, col sep=comma] {data/k_up/k_up_deg_3_cu_comparison.csv};
\legend{Avg Time, Max Time}
\end{axis}
\end{tikzpicture}
\caption{Random bounded degree trees with uniform costs - $O(\log \log n)$-approximation algorithm: runtime}
\label{fig:bounded_degree_c_uniform_kup_time}
\end{figure}

\begin{figure}[ht]
\centering
\begin{tikzpicture}
\begin{axis}[
    xlabel={Tree size ($n$)},
    ylabel={Approximation ratio},
    width=\textwidth,
    height=7cm,
    grid=major,
    legend pos=north west,
    mark size=2pt
]
\addplot[blue, mark=*, thick] table[x=Size, y=AvgRatio, col sep=comma] {data/k_up/k_up_deg_3_cu_comparison.csv};
\addplot[red, mark=square*, thick] table[x=Size, y=MaxRatio, col sep=comma] {data/k_up/k_up_deg_3_cu_comparison.csv};
\legend{Avg Ratio, Max Ratio}
\end{axis}
\end{tikzpicture}
\caption{Random bounded degree trees with uniform costs - $O(\log \log n)$-approximation algorithm: approximation quality}
\label{fig:bounded_degree_c_uniform_kup}
\end{figure}
\FloatBarrier
We observe that:
\begin{itemize}
\item The running time of the exact procedure almost strictly follows a linear function on the logarithmic plot. This will also be the case for the other cost distrubtions for bounded degree instances.
\item The $O\br{\log\log n}$-approximation algorithm only managed to process instances up to 12 vertices. We have observed that for the first instance having size $n=13$ which the algorithm encountered the processing was not completed despite few days of computation. This suggests that the running time of the  algorithm is unstable and largely dependent on the input structure.
\end{itemize}


\input{chapters/experimental_results/trees_worst_case_non_uniform_costs/bounded_degree_c_binomial.tex}
\clearpage\subsection{Random bounded degree trees with Exponential Cost Distribution ($T_{\Delta} || V, c \sim \mathcal{E} || \COST_{max}$)}
\label{sec:bounded_degree_c_exponential}

\subsubsection{Exact Algorithm}

\begin{table}[ht]
\centering
\tiny
\caption{Random bounded degree trees ($\Delta(T) = 3$) with exponential cost distribution - exact algorithm}
\label{tab:bounded_degree_c_exponential_dp}
\begin{tabular}{|c|c|c|}
\hline
$n$ & Avg Time (s) & Max Time (s) \\
\hline
\csvreader[head to column names, late after line=\\]{data/dp_tree/dp_tree_deg_3_ce_analysis.csv}{Size=\Size,AvgCPUTime_s=\AvgTime,MaxCPUTime_s=\MaxTime,AvgCost=\AvgCost,MaxCost=\MaxCost,InstanceCount=\Instances}{%
\Size & \AvgTime & \MaxTime}
\hline
\end{tabular}
\end{table}

\begin{figure}[ht]
\centering
\begin{tikzpicture}
\begin{axis}[
    xlabel={Tree size ($n$)},
    ylabel={Runtime (s)},
    ymode=log,
    width=\textwidth,
    height=7cm,
    grid=major,
    legend pos=north west,
    mark size=2pt
]
\addplot[blue, mark=*, thick] table[x=Size, y=AvgCPUTime_s, col sep=comma] {data/dp_tree/dp_tree_deg_3_ce_analysis.csv};
\addplot[red, mark=square*, thick] table[x=Size, y=MaxCPUTime_s, col sep=comma] {data/dp_tree/dp_tree_deg_3_ce_analysis.csv};
\legend{Avg Time, Max Time}
\end{axis}
\end{tikzpicture}
\caption{Random bounded degree trees with exponential costs - exact algorithm: runtime comparison}
\label{fig:bounded_degree_c_exponential_dp_time}
\end{figure}

\FloatBarrier

\subsubsection{$O(\log n / \log \log n)$-Approximation Algorithm}

\begin{table}[ht]
\centering
\tiny
\caption{Random bounded degree trees ($\Delta(T) = 3$) with exponential cost distribution - $O(\log n / \log \log n)$-approximation algorithm}
\label{tab:bounded_degree_c_exponential_cic}
\begin{tabular}{|c|c|c|c|c|}
\hline
$n$ & Avg Time (s) & Max Time (s) & Avg Ratio & Max Ratio \\
\hline
\csvreader[head to column names, late after line=\\]{data/cic_insp/cic_insp_deg_3_ce_comparison.csv}{Size=\Size,AvgCPUTime_s=\AvgTime,MaxCPUTime_s=\MaxTime,AvgCost=\AvgCost,MaxCost=\MaxCost,AvgRatio=\AvgRatio,MaxRatio=\MaxRatio}{%
\Size & \AvgTime & \MaxTime & \AvgRatio & \MaxRatio}
\hline
\end{tabular}
\end{table}

\begin{figure}[ht]
\centering
\begin{tikzpicture}
\begin{axis}[
    xlabel={Tree size ($n$)},
    ylabel={Runtime (s)},
    ymode=log,
    width=\textwidth,
    height=7cm,
    grid=major,
    legend pos=north west,
    mark size=2pt
]
\addplot[blue, mark=*, thick] table[x=Size, y=AvgCPUTime_s, col sep=comma] {data/cic_insp/cic_insp_deg_3_ce_comparison.csv};
\addplot[red, mark=square*, thick] table[x=Size, y=MaxCPUTime_s, col sep=comma] {data/cic_insp/cic_insp_deg_3_ce_comparison.csv};
\legend{Avg Time, Max Time}
\end{axis}
\end{tikzpicture}
\caption{Random bounded degree trees with exponential costs - $O(\log n / \log \log n)$-approximation algorithm: runtime}
\label{fig:bounded_degree_c_exponential_cic_time}
\end{figure}

\begin{figure}[ht]
\centering
\begin{tikzpicture}
\begin{axis}[
    xlabel={Tree size ($n$)},
    ylabel={Approximation ratio},
    width=\textwidth,
    height=7cm,
    grid=major,
    legend pos=north west,
    mark size=2pt
]
\addplot[blue, mark=*, thick] table[x=Size, y=AvgRatio, col sep=comma] {data/cic_insp/cic_insp_deg_3_ce_comparison.csv};
\addplot[red, mark=square*, thick] table[x=Size, y=MaxRatio, col sep=comma] {data/cic_insp/cic_insp_deg_3_ce_comparison.csv};
\legend{Avg Ratio, Max Ratio}
\end{axis}
\end{tikzpicture}
\caption{Random bounded degree trees with exponential costs - $O(\log n / \log \log n)$-approximation algorithm: approximation quality}
\label{fig:bounded_degree_c_exponential_cic}
\end{figure}

\FloatBarrier

\subsubsection{$O(\sqrt{\log n})$-Approximation Algorithm}

\begin{table}[ht]
\centering
\tiny
\caption{Random bounded degree trees ($\Delta(T) = 3$) with exponential cost distribution - $O(\sqrt{\log n})$-approximation algorithm}
\label{tab:bounded_degree_c_exponential_der}
\begin{tabular}{|c|c|c|c|c|}
\hline
$n$ & Avg Time (s) & Max Time (s) & Avg Ratio & Max Ratio \\
\hline
\csvreader[head to column names, late after line=\\]{data/der_insp/der_insp_deg_3_ce_comparison.csv}{Size=\Size,AvgCPUTime_s=\AvgTime,MaxCPUTime_s=\MaxTime,AvgCost=\AvgCost,MaxCost=\MaxCost,AvgRatio=\AvgRatio,MaxRatio=\MaxRatio}{%
\Size & \AvgTime & \MaxTime & \AvgRatio & \MaxRatio}
\hline
\end{tabular}
\end{table}

\begin{figure}[ht]
\centering
\begin{tikzpicture}
\begin{axis}[
    xlabel={Tree size ($n$)},
    ylabel={Runtime (s)},
    ymode=log,
    width=\textwidth,
    height=7cm,
    grid=major,
    legend pos=north west,
    mark size=2pt
]
\addplot[blue, mark=*, thick] table[x=Size, y=AvgCPUTime_s, col sep=comma] {data/der_insp/der_insp_deg_3_ce_comparison.csv};
\addplot[red, mark=square*, thick] table[x=Size, y=MaxCPUTime_s, col sep=comma] {data/der_insp/der_insp_deg_3_ce_comparison.csv};
\legend{Avg Time, Max Time}
\end{axis}
\end{tikzpicture}
\caption{Random bounded degree trees with exponential costs - $O(\sqrt{\log n})$-approximation algorithm: runtime}
\label{fig:bounded_degree_c_exponential_der_time}
\end{figure}

\begin{figure}[ht]
\centering
\begin{tikzpicture}
\begin{axis}[
    xlabel={Tree size ($n$)},
    ylabel={Approximation ratio},
    width=\textwidth,
    height=7cm,
    grid=major,
    legend pos=north west,
    mark size=2pt
]
\addplot[blue, mark=*, thick] table[x=Size, y=AvgRatio, col sep=comma] {data/der_insp/der_insp_deg_3_ce_comparison.csv};
\addplot[red, mark=square*, thick] table[x=Size, y=MaxRatio, col sep=comma] {data/der_insp/der_insp_deg_3_ce_comparison.csv};
\legend{Avg Ratio, Max Ratio}
\end{axis}
\end{tikzpicture}
\caption{Random bounded degree trees with exponential costs - $O(\sqrt{\log n})$-approximation algorithm: approximation quality}
\label{fig:bounded_degree_c_exponential_der}
\end{figure}

\FloatBarrier

\subsubsection{$O(\log \log n)$-Approximation Algorithm}

\begin{table}[ht]
\centering
\tiny
\caption{Random bounded degree trees ($\Delta(T) = 3$) with exponential cost distribution - $O(\log \log n)$-approximation algorithm}
\label{tab:bounded_degree_c_exponential_kup}
\begin{tabular}{|c|c|c|c|c|}
\hline
$n$ & Avg Time (s) & Max Time (s) & Avg Ratio & Max Ratio \\
\hline
\csvreader[head to column names, late after line=\\]{data/k_up/k_up_deg_3_ce_comparison.csv}{Size=\Size,AvgCPUTime_s=\AvgTime,MaxCPUTime_s=\MaxTime,AvgCost=\AvgCost,MaxCost=\MaxCost,AvgRatio=\AvgRatio,MaxRatio=\MaxRatio}{%
\Size & \AvgTime & \MaxTime & \AvgRatio & \MaxRatio}
\hline
\end{tabular}
\end{table}

\begin{figure}[ht]
\centering
\begin{tikzpicture}
\begin{axis}[
    xlabel={Tree size ($n$)},
    ylabel={Runtime (s)},
    ymode=log,
    width=\textwidth,
    height=7cm,
    grid=major,
    legend pos=north west,
    mark size=2pt
]
\addplot[blue, mark=*, thick] table[x=Size, y=AvgCPUTime_s, col sep=comma] {data/k_up/k_up_deg_3_ce_comparison.csv};
\addplot[red, mark=square*, thick] table[x=Size, y=MaxCPUTime_s, col sep=comma] {data/k_up/k_up_deg_3_ce_comparison.csv};
\legend{Avg Time, Max Time}
\end{axis}
\end{tikzpicture}
\caption{Random bounded degree trees with exponential costs - $O(\log \log n)$-approximation algorithm: runtime}
\label{fig:bounded_degree_c_exponential_kup_time}
\end{figure}

\begin{figure}[ht]
\centering
\begin{tikzpicture}
\begin{axis}[
    xlabel={Tree size ($n$)},
    ylabel={Approximation ratio},
    width=\textwidth,
    height=7cm,
    grid=major,
    legend pos=north west,
    mark size=2pt
]
\addplot[blue, mark=*, thick] table[x=Size, y=AvgRatio, col sep=comma] {data/k_up/k_up_deg_3_ce_comparison.csv};
\addplot[red, mark=square*, thick] table[x=Size, y=MaxRatio, col sep=comma] {data/k_up/k_up_deg_3_ce_comparison.csv};
\legend{Avg Ratio, Max Ratio}
\end{axis}
\end{tikzpicture}
\caption{Random bounded degree trees with exponential costs - $O(\log \log n)$-approximation algorithm: approximation quality}
\label{fig:bounded_degree_c_exponential_kup}
\end{figure}
\FloatBarrier
This time the $O\br{\sqrt{\log n}}$-approximation algorithm exhibits the best approximation factors never exceeding the value of 2 as opposed to the two other algorithms.



\clearpage\subsection{Random bounded diameter trees with Uniform Cost Distribution ($T_{D} || V, c \sim \mathcal{U} || \COST_{max}$)}
\label{sec:bounded_diameter_c_uniform}

\subsubsection{Exact Algorithm}

\begin{table}[ht]
\centering
\tiny
\caption{Random bounded diameter trees ($\diam\br{T} = 6$) with uniform cost distribution - exact algorithm}
\label{tab:bounded_diameter_c_uniform_dp}
\begin{tabular}{|c|c|c|}
\hline
$n$ & Avg Time (s) & Max Time (s) \\
\hline
\csvreader[head to column names, late after line=\\]{data/dp_tree/dp_tree_diam_6_cu_analysis.csv}{Size=\Size,AvgCPUTime_s=\AvgTime,MaxCPUTime_s=\MaxTime,AvgCost=\AvgCost,MaxCost=\MaxCost,InstanceCount=\Instances}{%
\Size & \AvgTime & \MaxTime}
\hline
\end{tabular}
\end{table}

\begin{figure}[ht]
\centering
\begin{tikzpicture}
\begin{axis}[
    xlabel={Tree size ($n$)},
    ylabel={Runtime (s)},
    ymode=log,
    width=\textwidth,
    height=7cm,
    grid=major,
    legend pos=north west,
    mark size=2pt
]
\addplot[blue, mark=*, thick] table[x=Size, y=AvgCPUTime_s, col sep=comma] {data/dp_tree/dp_tree_diam_6_cu_analysis.csv};
\addplot[red, mark=square*, thick] table[x=Size, y=MaxCPUTime_s, col sep=comma] {data/dp_tree/dp_tree_diam_6_cu_analysis.csv};
\legend{Avg Time, Max Time}
\end{axis}
\end{tikzpicture}
\caption{Random bounded diameter trees with uniform costs - exact algorithm: runtime comparison}
\label{fig:bounded_diameter_c_uniform_dp_time}
\end{figure}

\FloatBarrier

\subsubsection{$O(\log n / \log \log n)$-Approximation Algorithm}

\begin{table}[ht]
\centering
\tiny
\caption{Random bounded diameter trees ($\diam\br{T} = 6$) with uniform cost distribution - $O(\log n / \log \log n)$-approximation algorithm}
\label{tab:bounded_diameter_c_uniform_cic}
\begin{tabular}{|c|c|c|c|c|}
\hline
$n$ & Avg Time (s) & Max Time (s) & Avg Ratio & Max Ratio \\
\hline
\csvreader[head to column names, late after line=\\]{data/cic_insp/cic_insp_diam_6_cu_comparison.csv}{Size=\Size,AvgCPUTime_s=\AvgTime,MaxCPUTime_s=\MaxTime,AvgCost=\AvgCost,MaxCost=\MaxCost,AvgRatio=\AvgRatio,MaxRatio=\MaxRatio}{%
\Size & \AvgTime & \MaxTime & \AvgRatio & \MaxRatio}
\hline
\end{tabular}
\end{table}

\begin{figure}[ht]
\centering
\begin{tikzpicture}
\begin{axis}[
    xlabel={Tree size ($n$)},
    ylabel={Runtime (s)},
    ymode=log,
    width=\textwidth,
    height=7cm,
    grid=major,
    legend pos=north west,
    mark size=2pt
]
\addplot[blue, mark=*, thick] table[x=Size, y=AvgCPUTime_s, col sep=comma] {data/cic_insp/cic_insp_diam_6_cu_comparison.csv};
\addplot[red, mark=square*, thick] table[x=Size, y=MaxCPUTime_s, col sep=comma] {data/cic_insp/cic_insp_diam_6_cu_comparison.csv};
\legend{Avg Time, Max Time}
\end{axis}
\end{tikzpicture}
\caption{Random bounded diameter trees with uniform costs - $O(\log n / \log \log n)$-approximation algorithm: runtime}
\label{fig:bounded_diameter_c_uniform_cic_time}
\end{figure}

\begin{figure}[ht]
\centering
\begin{tikzpicture}
\begin{axis}[
    xlabel={Tree size ($n$)},
    ylabel={Approximation ratio},
    width=\textwidth,
    height=7cm,
    grid=major,
    legend pos=north west,
    mark size=2pt
]
\addplot[blue, mark=*, thick] table[x=Size, y=AvgRatio, col sep=comma] {data/cic_insp/cic_insp_diam_6_cu_comparison.csv};
\addplot[red, mark=square*, thick] table[x=Size, y=MaxRatio, col sep=comma] {data/cic_insp/cic_insp_diam_6_cu_comparison.csv};
\legend{Avg Ratio, Max Ratio}
\end{axis}
\end{tikzpicture}
\caption{Random bounded diameter trees with uniform costs - $O(\log n / \log \log n)$-approximation algorithm: approximation quality}
\label{fig:bounded_diameter_c_uniform_cic}
\end{figure}

\FloatBarrier

\subsubsection{$O(\sqrt{\log n})$-Approximation Algorithm}

\begin{table}[ht]
\centering
\tiny
\caption{Random bounded diameter trees ($\diam\br{T} = 6$) with uniform cost distribution - $O(\sqrt{\log n})$-approximation algorithm}
\label{tab:bounded_diameter_c_uniform_der}
\begin{tabular}{|c|c|c|c|c|}
\hline
$n$ & Avg Time (s) & Max Time (s) & Avg Ratio & Max Ratio \\
\hline
\csvreader[head to column names, late after line=\\]{data/der_insp/der_insp_diam_6_cu_comparison.csv}{Size=\Size,AvgCPUTime_s=\AvgTime,MaxCPUTime_s=\MaxTime,AvgCost=\AvgCost,MaxCost=\MaxCost,AvgRatio=\AvgRatio,MaxRatio=\MaxRatio}{%
\Size & \AvgTime & \MaxTime & \AvgRatio & \MaxRatio}
\hline
\end{tabular}
\end{table}

\begin{figure}[ht]
\centering
\begin{tikzpicture}
\begin{axis}[
    xlabel={Tree size ($n$)},
    ylabel={Runtime (s)},
    ymode=log,
    width=\textwidth,
    height=7cm,
    grid=major,
    legend pos=north west,
    mark size=2pt
]
\addplot[blue, mark=*, thick] table[x=Size, y=AvgCPUTime_s, col sep=comma] {data/der_insp/der_insp_diam_6_cu_comparison.csv};
\addplot[red, mark=square*, thick] table[x=Size, y=MaxCPUTime_s, col sep=comma] {data/der_insp/der_insp_diam_6_cu_comparison.csv};
\legend{Avg Time, Max Time}
\end{axis}
\end{tikzpicture}
\caption{Random bounded diameter trees with uniform costs - $O(\sqrt{\log n})$-approximation algorithm: runtime}
\label{fig:bounded_diameter_c_uniform_der_time}
\end{figure}

\begin{figure}[ht]
\centering
\begin{tikzpicture}
\begin{axis}[
    xlabel={Tree size ($n$)},
    ylabel={Approximation ratio},
    width=\textwidth,
    height=7cm,
    grid=major,
    legend pos=north west,
    mark size=2pt
]
\addplot[blue, mark=*, thick] table[x=Size, y=AvgRatio, col sep=comma] {data/der_insp/der_insp_diam_6_cu_comparison.csv};
\addplot[red, mark=square*, thick] table[x=Size, y=MaxRatio, col sep=comma] {data/der_insp/der_insp_diam_6_cu_comparison.csv};
\legend{Avg Ratio, Max Ratio}
\end{axis}
\end{tikzpicture}
\caption{Random bounded diameter trees with uniform costs - $O(\sqrt{\log n})$-approximation algorithm: approximation quality}
\label{fig:bounded_diameter_c_uniform_der}
\end{figure}

\FloatBarrier

\subsubsection{$O(\log \log n)$-Approximation Algorithm}

\begin{table}[ht]
\centering
\tiny
\caption{Random bounded diameter trees ($\diam\br{T} = 6$) with uniform cost distribution - $O(\log \log n)$-approximation algorithm}
\label{tab:bounded_diameter_c_uniform_kup}
\begin{tabular}{|c|c|c|c|c|}
\hline
$n$ & Avg Time (s) & Max Time (s) & Avg Ratio & Max Ratio \\
\hline
\csvreader[head to column names, late after line=\\]{data/k_up/k_up_diam_6_cu_comparison.csv}{Size=\Size,AvgCPUTime_s=\AvgTime,MaxCPUTime_s=\MaxTime,AvgCost=\AvgCost,MaxCost=\MaxCost,AvgRatio=\AvgRatio,MaxRatio=\MaxRatio}{%
\Size & \AvgTime & \MaxTime & \AvgRatio & \MaxRatio}
\hline
\end{tabular}
\end{table}

\begin{figure}[ht]
\centering
\begin{tikzpicture}
\begin{axis}[
    xlabel={Tree size ($n$)},
    ylabel={Runtime (s)},
    ymode=log,
    width=\textwidth,
    height=7cm,
    grid=major,
    legend pos=north west,
    mark size=2pt
]
\addplot[blue, mark=*, thick] table[x=Size, y=AvgCPUTime_s, col sep=comma] {data/k_up/k_up_diam_6_cu_comparison.csv};
\addplot[red, mark=square*, thick] table[x=Size, y=MaxCPUTime_s, col sep=comma] {data/k_up/k_up_diam_6_cu_comparison.csv};
\legend{Avg Time, Max Time}
\end{axis}
\end{tikzpicture}
\caption{Random bounded diameter trees with uniform costs - $O(\log \log n)$-approximation algorithm: runtime}
\label{fig:bounded_diameter_c_uniform_kup_time}
\end{figure}

\begin{figure}[ht]
\centering
\begin{tikzpicture}
\begin{axis}[
    xlabel={Tree size ($n$)},
    ylabel={Approximation ratio},
    width=\textwidth,
    height=7cm,
    grid=major,
    legend pos=north west,
    mark size=2pt
]
\addplot[blue, mark=*, thick] table[x=Size, y=AvgRatio, col sep=comma] {data/k_up/k_up_diam_6_cu_comparison.csv};
\addplot[red, mark=square*, thick] table[x=Size, y=MaxRatio, col sep=comma] {data/k_up/k_up_diam_6_cu_comparison.csv};
\legend{Avg Ratio, Max Ratio}
\end{axis}
\end{tikzpicture}
\caption{Random bounded diameter trees with uniform costs - $O(\log \log n)$-approximation algorithm: approximation quality}
\label{fig:bounded_diameter_c_uniform_kup}
\end{figure}
\FloatBarrier
We see that all of the algorithms achieve the average approximation ratio of almost 1. This will usually also be the case for the other cost distributions for bounded diameter. Notably the approximation ratio of the $O\br{\log n/\log \log n}$-approximation algorithm never exceeds 1.6.

\input{chapters/experimental_results/trees_worst_case_non_uniform_costs/bounded_diameter_c_binomial.tex}
\input{chapters/experimental_results/trees_worst_case_non_uniform_costs/bounded_diameter_c_exponential.tex}


% \newpage
\section{Trees with Both Non-Uniform Weights and Costs (Average Case)}

This section presents experimental results for trees with both non-uniform vertex weights and edge costs in the average-case setting.

\input{chapters/experimental_results/non_uniform_weights_costs_average/random_wc.tex}
\input{chapters/experimental_results/non_uniform_weights_costs_average/bounded_degree_wc.tex}
\input{chapters/experimental_results/non_uniform_weights_costs_average/bounded_diameter_wc.tex}
\subsection{Large Degree Trees with Non-Uniform Weights and Costs ($T_{\delta} || V, w, c || \sum w_i C_i$)}
\label{sec:large_degree_wc}

Results for large degree trees with both non-uniform vertex weights and query costs (average case).

\begin{table}[H]
\centering
\caption{Large Degree trees with non-uniform weights and costs (average case)}
\label{tab:large_degree_wc}
\begin{tabular}{|c|c|c|c|c|}
\hline
$n$ & Avg Cost & Approx Ratio & Avg Time (ms) & Std Dev \\
\hline
\csvreader[head to column names]{data/large_degree_wc.csv}{}{%
\n & \avgcost & \approxratio & \avgtime & \stddev \\
}
\hline
\end{tabular}
\end{table}
\begin{figure}[H]
\centering
\begin{tikzpicture}
\begin{axis}[
    xlabel={Tree size ($n$)},
    ylabel={Average cost},
    width=0.48\textwidth,
    height=6cm,
    grid=major,
    legend pos=north west,
    mark size=2pt
]
\addplot[blue, mark=*, thick] table[x=n, y=avgcost, col sep=comma] {data/large_degree_wc.csv};
\legend{Avg Cost}
\end{axis}
\end{tikzpicture}
\hfill
\begin{tikzpicture}
\begin{axis}[
    xlabel={Tree size ($n$)},
    ylabel={Runtime (ms)},
    width=0.48\textwidth,
    height=6cm,
    grid=major,
    legend pos=north west,
    mark size=2pt
]
\addplot[red, mark=square*, thick] table[x=n, y=avgtime, col sep=comma] {data/large_degree_wc.csv};
\legend{Avg Time}
\end{axis}
\end{tikzpicture}
\caption{Performance metrics: Large Degree trees with non-uniform weights and costs}
\label{fig:large_degree_wc}
\end{figure}


\input{chapters/experimental_results/non_uniform_weights_costs_average/large_diameter_wc.tex}
\input{chapters/experimental_results/non_uniform_weights_costs_average/star_wc.tex}
\input{chapters/experimental_results/non_uniform_weights_costs_average/spider_wc.tex}
\subsection{Caterpillar Trees with Non-Uniform Weights and Costs ($T_{cat} || V, w, c || \sum w_i C_i$)}
\label{sec:caterpillar_wc}

Results for caterpillar trees with both non-uniform vertex weights and query costs (average case).

\begin{table}[H]
\centering
\caption{Caterpillar trees with non-uniform weights and costs (average case)}
\label{tab:caterpillar_wc}
\begin{tabular}{|c|c|c|c|c|}
\hline
$n$ & Avg Cost & Approx Ratio & Avg Time (ms) & Std Dev \\
\hline
\csvreader[head to column names]{data/caterpillar_wc.csv}{}{%
\n & \avgcost & \approxratio & \avgtime & \stddev \\
}
\hline
\end{tabular}
\end{table}
\begin{figure}[H]
\centering
\begin{tikzpicture}
\begin{axis}[
    xlabel={Tree size ($n$)},
    ylabel={Average cost},
    width=0.48\textwidth,
    height=6cm,
    grid=major,
    legend pos=north west,
    mark size=2pt
]
\addplot[blue, mark=*, thick] table[x=n, y=avgcost, col sep=comma] {data/caterpillar_wc.csv};
\legend{Avg Cost}
\end{axis}
\end{tikzpicture}
\hfill
\begin{tikzpicture}
\begin{axis}[
    xlabel={Tree size ($n$)},
    ylabel={Runtime (ms)},
    width=0.48\textwidth,
    height=6cm,
    grid=major,
    legend pos=north west,
    mark size=2pt
]
\addplot[red, mark=square*, thick] table[x=n, y=avgtime, col sep=comma] {data/caterpillar_wc.csv};
\legend{Avg Time}
\end{axis}
\end{tikzpicture}
\caption{Performance metrics: Caterpillar trees with non-uniform weights and costs}
\label{fig:caterpillar_wc}
\end{figure}


\input{chapters/experimental_results/non_uniform_weights_costs_average/balanced_wc.tex}
\input{chapters/experimental_results/non_uniform_weights_costs_average/broom_wc.tex}
\subsection{Lobster Trees with Non-Uniform Weights and Costs ($T_{lob} || V, w, c || \sum w_i C_i$)}
\label{sec:lobster_wc}

Results for lobster trees with both non-uniform vertex weights and query costs (average case).

\begin{table}[H]
\centering
\caption{Lobster trees with non-uniform weights and costs (average case)}
\label{tab:lobster_wc}
\begin{tabular}{|c|c|c|c|c|}
\hline
$n$ & Avg Cost & Approx Ratio & Avg Time (ms) & Std Dev \\
\hline
\csvreader[head to column names]{data/lobster_wc.csv}{}{%
\n & \avgcost & \approxratio & \avgtime & \stddev \\
}
\hline
\end{tabular}
\end{table}
\begin{figure}[H]
\centering
\begin{tikzpicture}
\begin{axis}[
    xlabel={Tree size ($n$)},
    ylabel={Average cost},
    width=0.48\textwidth,
    height=6cm,
    grid=major,
    legend pos=north west,
    mark size=2pt
]
\addplot[blue, mark=*, thick] table[x=n, y=avgcost, col sep=comma] {data/lobster_wc.csv};
\legend{Avg Cost}
\end{axis}
\end{tikzpicture}
\hfill
\begin{tikzpicture}
\begin{axis}[
    xlabel={Tree size ($n$)},
    ylabel={Runtime (ms)},
    width=0.48\textwidth,
    height=6cm,
    grid=major,
    legend pos=north west,
    mark size=2pt
]
\addplot[red, mark=square*, thick] table[x=n, y=avgtime, col sep=comma] {data/lobster_wc.csv};
\legend{Avg Time}
\end{axis}
\end{tikzpicture}
\caption{Performance metrics: Lobster trees with non-uniform weights and costs}
\label{fig:lobster_wc}
\end{figure}


\subsection{Full m-ary Trees with Non-Uniform Weights and Costs ($T_m || V, w, c || \sum w_i C_i$)}
\label{sec:mary_wc}

Results for full m-ary trees with both non-uniform vertex weights and query costs (average case).

\begin{table}[H]
\centering
\caption{Full m-ary trees with non-uniform weights and costs (average case)}
\label{tab:mary_wc}
\begin{tabular}{|c|c|c|c|c|}
\hline
$n$ & Avg Cost & Approx Ratio & Avg Time (ms) & Std Dev \\
\hline
\csvreader[head to column names]{data/mary_wc.csv}{}{%
\n & \avgcost & \approxratio & \avgtime & \stddev \\
}
\hline
\end{tabular}
\end{table}
\begin{figure}[H]
\centering
\begin{tikzpicture}
\begin{axis}[
    xlabel={Tree size ($n$)},
    ylabel={Average cost},
    width=0.48\textwidth,
    height=6cm,
    grid=major,
    legend pos=north west,
    mark size=2pt
]
\addplot[blue, mark=*, thick] table[x=n, y=avgcost, col sep=comma] {data/mary_wc.csv};
\legend{Avg Cost}
\end{axis}
\end{tikzpicture}
\hfill
\begin{tikzpicture}
\begin{axis}[
    xlabel={Tree size ($n$)},
    ylabel={Runtime (ms)},
    width=0.48\textwidth,
    height=6cm,
    grid=major,
    legend pos=north west,
    mark size=2pt
]
\addplot[red, mark=square*, thick] table[x=n, y=avgtime, col sep=comma] {data/mary_wc.csv};
\legend{Avg Time}
\end{axis}
\end{tikzpicture}
\caption{Performance metrics: Full m-ary trees with non-uniform weights and costs}
\label{fig:mary_wc}
\end{figure}


\input{chapters/experimental_results/non_uniform_weights_costs_average/probabilistic_wc.tex}
\subsection{Recursive Trees with Non-Uniform Weights and Costs ($T_{rec} || V, w, c || \sum w_i C_i$)}
\label{sec:recursive_wc}

Results for recursive trees with both non-uniform vertex weights and query costs (average case).

\begin{table}[H]
\centering
\caption{Recursive trees with non-uniform weights and costs (average case)}
\label{tab:recursive_wc}
\begin{tabular}{|c|c|c|c|c|}
\hline
$n$ & Avg Cost & Approx Ratio & Avg Time (ms) & Std Dev \\
\hline
\csvreader[head to column names]{data/recursive_wc.csv}{}{%
\n & \avgcost & \approxratio & \avgtime & \stddev \\
}
\hline
\end{tabular}
\end{table}
\begin{figure}[H]
\centering
\begin{tikzpicture}
\begin{axis}[
    xlabel={Tree size ($n$)},
    ylabel={Average cost},
    width=0.48\textwidth,
    height=6cm,
    grid=major,
    legend pos=north west,
    mark size=2pt
]
\addplot[blue, mark=*, thick] table[x=n, y=avgcost, col sep=comma] {data/recursive_wc.csv};
\legend{Avg Cost}
\end{axis}
\end{tikzpicture}
\hfill
\begin{tikzpicture}
\begin{axis}[
    xlabel={Tree size ($n$)},
    ylabel={Runtime (ms)},
    width=0.48\textwidth,
    height=6cm,
    grid=major,
    legend pos=north west,
    mark size=2pt
]
\addplot[red, mark=square*, thick] table[x=n, y=avgtime, col sep=comma] {data/recursive_wc.csv};
\legend{Avg Time}
\end{axis}
\end{tikzpicture}
\caption{Performance metrics: Recursive trees with non-uniform weights and costs}
\label{fig:recursive_wc}
\end{figure}


\subsection{Preferential Attachment Trees with Non-Uniform Weights and Costs ($T_{PA} || V, w, c || \sum w_i C_i$)}
\label{sec:preferential_wc}

Results for preferential attachment trees with both non-uniform vertex weights and query costs (average case).

\begin{table}[H]
\centering
\caption{Preferential Attachment trees with non-uniform weights and costs (average case)}
\label{tab:preferential_wc}
\begin{tabular}{|c|c|c|c|c|}
\hline
$n$ & Avg Cost & Approx Ratio & Avg Time (ms) & Std Dev \\
\hline
\csvreader[head to column names]{data/preferential_wc.csv}{}{%
\n & \avgcost & \approxratio & \avgtime & \stddev \\
}
\hline
\end{tabular}
\end{table}
\begin{figure}[H]
\centering
\begin{tikzpicture}
\begin{axis}[
    xlabel={Tree size ($n$)},
    ylabel={Average cost},
    width=0.48\textwidth,
    height=6cm,
    grid=major,
    legend pos=north west,
    mark size=2pt
]
\addplot[blue, mark=*, thick] table[x=n, y=avgcost, col sep=comma] {data/preferential_wc.csv};
\legend{Avg Cost}
\end{axis}
\end{tikzpicture}
\hfill
\begin{tikzpicture}
\begin{axis}[
    xlabel={Tree size ($n$)},
    ylabel={Runtime (ms)},
    width=0.48\textwidth,
    height=6cm,
    grid=major,
    legend pos=north west,
    mark size=2pt
]
\addplot[red, mark=square*, thick] table[x=n, y=avgtime, col sep=comma] {data/preferential_wc.csv};
\legend{Avg Time}
\end{axis}
\end{tikzpicture}
\caption{Performance metrics: Preferential Attachment trees with non-uniform weights and costs}
\label{fig:preferential_wc}
\end{figure}




